%% =====================================================================
%%  Capítulo 12 -- Linealización por retroalimentación
%%  Formato nuevo (noviembre de 2026). Integra: Handout9_2.tex (versión
%%  anterior), las ecuaciones de las dos clases de 2026
%%  (EcuacionesFL2026.tex, FeedbackLin1_2026 y FeedbackLin2_2026) y sus
%%  scripts FL1_V*/FL2_V*, reescritos en codigo/cap12_linealizacion.
%%  Correcciones respecto a las versiones anteriores: ver \S\ref{sec:fl-cierre}.
%%  Números verificados con herramientas/cap12_numeros.py (88/88) y
%%  codigo/cap12_linealizacion/verificarCap12.m (37/37 en Octave 8.4).
%% =====================================================================

\capitulo[cap:linealizacion]{Linealización por retroalimentación}
  {Cuanto más rico y profundo es el conocimiento parece que se vuelve más contumaz la ignorancia.}
  {Antonio Muñoz Molina}

\begin{objetivos}[partes/12-linealizacion.html]
  \item Diseñar un controlador sobre la linealización jacobiana, con y sin acción integral, y decir qué garantiza y qué no.
  \item Calcular derivadas de Lie, reconocer un difeomorfismo local (y uno que no es global) y escribir un sistema afín en coordenadas nuevas.
  \item Calcular el grado relativo de una salida, reconocer los puntos donde no está definido y llevar el sistema a la forma normal.
  \item Diseñar la ley que linealiza el mapa de la entrada a la salida, y un lazo externo lineal para estabilizar y para seguir referencias.
  \item Calcular la dinámica cero, decidir si un sistema es de fase mínima y concluir estabilidad local; saber qué más hace falta para un resultado global.
  \item Reconocer los límites de la cancelación: singularidades, error de modelo, no linealidades que ayudan.
\end{objetivos}

%% ---------------------------------------------------------------------
\section{Cancelar en vez de aproximar}
\label{sec:fl-por-que}
%% ---------------------------------------------------------------------

Hasta aquí, cada herramienta de diseño del curso ha sido una \emph{función de energía}: una función de Lyapunov con $\Vdot<0$, una función de Lyapunov de control, una ganancia de ISS que acota lo que entra, un almacenamiento que cumple $\Vdot\le u\T y$. Este capítulo cambia de filosofía. La pregunta ahora es geométrica: ¿existe un \clave{cambio de coordenadas} y una \clave{realimentación} que conviertan el sistema no lineal en uno \emph{lineal}, de manera exacta y no aproximada? Si existen, todo lo que se sabe de control lineal (ubicación de polos, seguimiento, LQR) queda disponible.

\notamargen{\textbf{Ruta de las clases.} Clase~1: \S\ref{sec:fl-jacobiana}--\ref{sec:fl-forma}. Clase~2: \S\ref{sec:fl-io}--\ref{sec:fl-estabilidad}. Las demostraciones de la \cref{sec:fl-grado} y la \cref{sec:fl-estabilidad} (lema de independencia y teorema local) se leen por cuenta propia.}

La respuesta tiene tres partes, y las tres son el tema del capítulo. Primera: para la clase de sistemas \emph{afines en la entrada} hay un procedimiento (el grado relativo y la forma normal) que dice exactamente qué se puede linealizar. Segunda: casi siempre queda una parte del sistema que la realimentación vuelve \emph{invisible} desde la salida, la \clave{dinámica cero}, y de ella depende la estabilidad. Tercera: cancelar exige conocer el modelo; la cancelación exacta puede ser frágil, puede pedir esfuerzos infinitos y puede destruir no linealidades que ayudaban.

\subsection{El problema de estabilización}

Para un sistema $\dot x=f(x,u)$ con salida $y=h(x)$, el \clave{problema de estabilización por realimentación de estado} es encontrar una ley \emph{estática} $u=\gamma(x)$ tal que un equilibrio deseado sea asintóticamente estable para el lazo cerrado $\dot x=f\bigl(x,\gamma(x)\bigr)$. Si el equilibrio deseado no es el origen, sino un punto $\eq$, hace falta antes un valor de régimen $u^*$ con
\begin{equation}
  0=f(\eq,u^*) .
  \label{eq:fl-regimen}
\end{equation}
Las tareas típicas de un controlador (estabilizar, seguir una referencia, rechazar perturbaciones, dar cierta respuesta transitoria, respetar límites de la señal de control) se plantean a partir de esta. El capítulo recorre tres caminos: (1) diseñar sobre la linealización jacobiana; (2) agregar acción integral para tolerar parámetros mal conocidos; (3) \emph{linealización por retroalimentación} (\emph{feedback linearization}): cancelar las no linealidades. En los tres se supone que el estado completo $x$ se mide.

\subsection{El ejemplo que nos acompaña: un levitador magnético}

Una esfera de material ferromagnético, de masa $m$, cuelga bajo un electroimán (\cref{fig:fl-levitador}). Con $y\ge0$ la distancia entre la esfera y el electroimán (medida hacia abajo), $k$ el coeficiente de fricción viscosa, $g$ la gravedad e $i$ la corriente de la bobina, la segunda ley de Newton da
\[
  m\ddot y=-k\dot y+mg+F(y,i) .
\]
La inductancia depende de la posición: es máxima con la esfera pegada y baja a $L_1$ cuando se aleja,
\begin{equation}
  L(y)=L_1+\frac{L_0}{1+y/a},
  \qquad
  E(y,i)=\tfrac12L(y)\,i^2,
  \qquad
  F(y,i)=\jac Ey=-\frac{L_0\,a\,i^2}{2(a+y)^2}<0 ,
  \label{eq:fl-fuerza}
\end{equation}
con $L_0,L_1,a>0$. El signo de $F$ dice algo físico que será decisivo en la \cref{sec:fl-levitador}: el electroimán \emph{solo atrae}. Con el voltaje $u$ como entrada y el flujo $\phi_m=L(y)\,i$, la ley de Kirchhoff es $u=\dot\phi_m+Ri$, y $\dot\phi_m=L'(y)\dot y\,i+L(y)\,\dot\imath$.

\begin{figure}[ht]
\centering
\includegraphics[width=0.5\linewidth]{cap12/t_levitador}
\caption{Levitador magnético: la corriente $i$ crea una fuerza $F(y,i)$ que \emph{atrae} la esfera; el voltaje $u$ es la entrada y la posición $y$ (medida hacia abajo) es la salida.}
\label{fig:fl-levitador}
\end{figure}

Con el estado $x=(x_1,x_2,x_3)=(y,\dot y,i)$, el modelo es
\begin{equation}
  \begin{aligned}
    \dot x_1&=x_2\\
    \dot x_2&=g-\frac km\,x_2-\frac{L_0\,a\,x_3^2}{2m(a+x_1)^2}\\
    \dot x_3&=\frac{1}{L(x_1)}\Bigl(u-Rx_3+\frac{L_0\,a\,x_2\,x_3}{(a+x_1)^2}\Bigr)
  \end{aligned}
  \qquad\qquad y=x_1 .
  \label{eq:fl-levitador}
\end{equation}
Usamos los valores de Khalil~\citep{Kha02}: $m=0.1$\,kg, $k=0.001$\,N\,s/m, $g=9.81$\,m/s$^2$, $a=0.05$\,m, $L_0=0.01$\,H, $L_1=0.02$\,H y $R=1$\,$\Omega$.

\notamargen{\textbf{Notación.} El grado relativo se escribe $r$ (Khalil usa $\rho$). La inductancia $L(y)$ siempre lleva su argumento; la derivada de Lie $L_fh$ siempre lleva el subíndice del campo. Las componentes de un cambio de coordenadas se escriben $\phi_i$, con subíndice; sin subíndice, $\phi$ sigue siendo la solución de la \cref{sec:pm-existencia}. La fricción del levitador es $k$; las ganancias de diseño, $k_i$.}

\begin{historia}[De Brockett a Kokotović]
En 1978, Roger Brockett preguntó qué propiedades de un sistema no lineal sobreviven a un cambio de coordenadas y una realimentación, y mostró que algunos sistemas se vuelven lineales con esas dos operaciones~\citep{Bro78}. Jakubczyk y Respondek (1980) dieron las condiciones exactas para que un sistema sea linealizable \emph{en todo su estado}~\citep{JaRe80}, y Hunt, Su y Meyer (1983) llevaron la idea al control de vuelo de helicópteros~\citep{HSM83}. Isidori, Krener, Gori-Giorgi y Monaco (1981) construyeron la versión geométrica para el problema de entrada y salida~\citep{IKGM81}. Byrnes e Isidori, entre 1984 y 1991, definieron la \emph{dinámica cero} y probaron los resultados de estabilización de los sistemas de \emph{fase mínima}~\citep{ByIs88,ByIs91}, y con Willems los conectaron con la pasividad~\citep{BIW91}. En 2001, Kokotović y Arcak hicieron el balance desde el diseño: la cancelación exacta puede ser frágil y puede desperdiciar no linealidades útiles~\citep{KoAr01}. El libro de Isidori~\citep{Isi95} es la referencia de toda la teoría.
\end{historia}

\begin{pausa}
Antes de seguir, apueste: si una ley de control hace que la salida $y(t)$ de un sistema converja a cero exponencialmente, con los polos que uno escoja, ¿converge a cero también el estado $x(t)$?
\respuesta{No necesariamente. La ley puede volver una parte del estado invisible desde la salida, y esa parte puede ser inestable: en el ejemplo de la \cref{sec:fl-cero}, $y\to0$ como se diseñó mientras el estado escapa a infinito en tiempo finito. Lo que decide es la \emph{dinámica cero} (\cref{def:fl-cero}).}
\end{pausa}

%% ---------------------------------------------------------------------
\section{Dos diseños sobre la linealización jacobiana}
\label{sec:fl-jacobiana}
%% ---------------------------------------------------------------------

\subsection{Equilibrios e inestabilidad}

Para sostener la esfera en $y=\bar y>0$ hace falta $\dot x=0$ con $x_1=\bar y$: de la primera ecuación de \eqref{eq:fl-levitador}, $x_2=0$; de la segunda, la fuerza magnética equilibra el peso; de la tercera, el voltaje solo cubre la caída en la resistencia:
\begin{equation}
  \eq(\bar y)=\Bigl(\bar y,\ 0,\ \sqrt{\tfrac{2gm}{L_0a}}\,(a+\bar y)\Bigr),
  \qquad
  u^*(\bar y)=R\,x_3^* .
  \label{eq:fl-equilibrio}
\end{equation}
Hay un equilibrio por cada $\bar y$: la familia está parametrizada por la referencia. Con $\bar y=0.05$\,m, $x_3^*=u^*=6.2642$ (amperios y voltios). La linealización jacobiana en $\eq(\bar y)$ es $\dot x_\delta=A(\bar y)\,x_\delta+B(\bar y)\,u_\delta$, con $x_\delta=x-\eq$, $u_\delta=u-u^*$ y, escribiendo $D=L_1(a+\bar y)+L_0a$,
\begin{equation}
  A(\bar y)=\begin{bmatrix}0&1&0\\[2pt] \dfrac{2g}{a+\bar y}&-\dfrac km&-\dfrac{\sqrt{2gaL_0/m}}{a+\bar y}\\[8pt] 0&\dfrac{\sqrt{2gmaL_0}}{D}&-\dfrac{R(a+\bar y)}{D}\end{bmatrix},
  \qquad
  B(\bar y)=\begin{bmatrix}0\\0\\ \dfrac{a+\bar y}{D}\end{bmatrix}.
  \label{eq:fl-AB}
\end{equation}
Con $\bar y=0.05$, los valores propios de $A$ son $13.641$, $-14.813$ y $-38.838$: por el método indirecto de Lyapunov, el equilibrio es \emph{inestable}. Y lo es para todo $\bar y$: la parte real del valor propio inestable baja con $\bar y$, pero nunca llega a cero (\cref{fig:fl-inestable}, izquierda). Note además que $A$ y $B$ \emph{dependen de} $\bar y$: una ganancia diseñada en un punto se aplica, en otro punto, a otra planta.

\begin{figure}[ht]
\centering
\includegraphics[width=\linewidth]{cap12/f_levitador_inestable}
\caption{Izquierda: partes reales de los valores propios de $A(\bar y)$; siempre hay uno positivo. Derecha: $b(x)=L_gL_f^2h(x)$ en función de la corriente para tres posiciones; se anula en $i=0$ (\cref{sec:fl-grado}). El punto es el equilibrio en $\bar y=0.05$.}
\label{fig:fl-inestable}
\end{figure}

\codigorepo{cap12_linealizacion/C12_Ej1_LevitadorJacobiano.m}{equilibrio}{Equilibrio y linealización jacobiana}

\subsection{Método 1: realimentación de estados sobre la jacobiana}

Si $(A,B)$ es controlable, se escoge $K$ con $A-BK$ Hurwitz y se aplica
\begin{equation}
  u=\underbrace{u^*}_{\text{prealimentación}}\ \underbrace{-\,K\,(x-\eq)}_{\text{realimentación}} .
  \label{eq:fl-metodo1}
\end{equation}
El lazo cerrado no lineal tiene en $\eq$ la linealización $\dot x_\delta=(A-BK)x_\delta$, y el método indirecto de Lyapunov garantiza que $\eq$ es localmente exponencialmente estable. Con los polos en $-20$, $-25$ y $-30$ para $\bar y=0.05$, $K=[-237.17,\ -16.013,\ 0.87475]$.

\begin{trampa}[label=tr:fl-local]{``Diseñé los polos; la respuesta es la que calculé''}
\plausible{Con $A-BK$ Hurwitz, la esfera responde como el modelo lineal $\dot x_\delta=(A-BK)x_\delta$.}
El método indirecto garantiza estabilidad \emph{local}: existe una vecindad de $\eq$ desde la que las soluciones convergen, pero el teorema no dice cuán grande es, ni cuánto se parece la respuesta a la lineal. Con el mismo $K$ y cuatro condiciones iniciales (\cref{fig:fl-jacobiana}): desde $y_0=0.045$ la diferencia con la predicción lineal es $2.8\times10^{-5}$\,m; desde $0.08$ y $0.01$, alrededor de $1$\,mm; desde $0.15$, $1.07$\,cm, con un voltaje de hasta $30$\,V.
\correcto{El diseño lineal \emph{predice} la respuesta solo cerca de $\eq$. Para saber desde dónde converge se estima la región de atracción con una función de Lyapunov: $V=x_\delta\T Px_\delta$ con $P(A-BK)+(A-BK)\T P=-Q$ y un conjunto de nivel de $V$ contenido en la zona donde $\Vdot<0$ para la planta \emph{no lineal}.}
\end{trampa}

\begin{figure}[ht]
\centering
\includegraphics[width=\linewidth]{cap12/f_jacobiana}
\caption{Mismo $K$ (diseñado en $\bar y=0.05$), cuatro condiciones iniciales con velocidad cero y la corriente de equilibrio. Izquierda: la planta (continua) y la predicción lineal (discontinua). Derecha: la diferencia crece con la distancia al punto de operación.}
\label{fig:fl-jacobiana}
\end{figure}

\codigorepo{cap12_linealizacion/C12_Ej1_LevitadorJacobiano.m}{comparacion}{La promesa lineal contra la planta}

\subsection{Método 2: control integral}

El método 1 necesita $u^*$, y $u^*$ sale del modelo. Si la masa real no es la del modelo, la prealimentación está mal y queda un error permanente. El remedio clásico es integrar el error de regulación. Para una planta con parámetros inciertos $w$, $\dot x=f(x,u,w)$, $y=h(x,w)$, y una referencia constante $\bar y$, se agrega el estado
\begin{equation}
  \dot\sigma=e=y-\bar y
  \label{eq:fl-integrador}
\end{equation}
y se estabiliza el sistema aumentado $(x,\sigma)$ con $u=-K_1x-K_2\sigma$ (más un valor nominal, si se quiere). En cualquier equilibrio del lazo cerrado, $\dot\sigma=0$ obliga a $y=\bar y$, \emph{sea cual sea} $w$; el valor $\sigma^*$ se acomoda solo para producir el $u^*$ que la planta real necesita.

\begin{proposicion}[label=prop:fl-integral]{Control integral vía linealización \citep[\S12.3]{Kha02}}
Sean $A=\jacat fx{(\eq,u^*)}$, $B=\jacat fu{(\eq,u^*)}$, $C=\jacat hx{\eq}$, con $p$ salidas reguladas. Si $(A,B)$ es controlable y
\begin{equation}
  \rank\begin{bmatrix}A&B\\ C&0\end{bmatrix}=n+p ,
  \label{eq:fl-rango}
\end{equation}
entonces el par aumentado $\Bigl(\begin{bmatrix}A&0\\ C&0\end{bmatrix},\begin{bmatrix}B\\0\end{bmatrix}\Bigr)$ es controlable, y existe $[K_1\ K_2]$ que ubica sus valores propios donde se quiera.
\end{proposicion}

\begin{proof}
Por el criterio de Hautus, hay que ver que $[\,A_a-\lambda I\ \ B_a\,]$ tiene rango $n+p$ para todo $\lambda\in\C$. Para $\lambda\neq0$, sus filas son $\begin{bmatrix}A-\lambda I&0&B\\ C&-\lambda I&0\end{bmatrix}$; el bloque $-\lambda I$ es invertible y $[\,A-\lambda I\ \ B\,]$ tiene rango $n$ por la controlabilidad de $(A,B)$. Para $\lambda=0$, la matriz es $\begin{bmatrix}A&0&B\\ C&0&0\end{bmatrix}$, cuyo rango es el de \eqref{eq:fl-rango}.
\end{proof}

En una planta de una entrada y una salida, la condición \eqref{eq:fl-rango} dice que la linealización no tiene un cero en $s=0$: un cero en el origen bloquearía exactamente la señal constante que el integrador necesita. En el levitador el rango es $4$, y con los polos aumentados en $-15$, $-20$, $-25$ y $-30$ se obtiene $K_1=[-482.15,\ -24.992,\ 1.2498]$, $K_2=-1795.9$. En $t=0.5$\,s la masa real pasa de $0.10$ a $0.12$\,kg y los dos controladores siguen usando el modelo nominal (\cref{fig:fl-integral}): la realimentación de estados se queda con un error de $1.03$\,cm; con el integrador el error se anula y el voltaje se reajusta a $6.8621$\,V, exactamente $R\,x_3^*$ calculado con $m=0.12$.

\begin{figure}[ht]
\centering
\includegraphics[width=\linewidth]{cap12/f_integral}
\caption{Una carga que el modelo no conoce: desde $t=0.5$\,s la masa es $0.12$\,kg. Izquierda: posición. Derecha: el voltaje; la línea punteada es el valor de equilibrio que exige la masa real.}
\label{fig:fl-integral}
\end{figure}

\codigorepo{cap12_linealizacion/C12_Ej2_ControlIntegral.m}{carga}{Realimentación de estados contra acción integral}

\begin{trampa}[label=tr:fl-integral]{``Con integrador, el controlador es robusto''}
\plausible{La acción integral elimina el error ante cualquier cambio de la planta.}
La acción integral elimina el error \emph{de régimen} ante perturbaciones y parámetros \emph{constantes}, siempre que el lazo cerrado siga teniendo un equilibrio estable cerca del nuevo punto. El diseño se hizo con la linealización en el punto nominal: un cambio grande de la planta (otra masa, otro punto de operación) puede volver inestable el lazo, y el integrador no lo impide. Tampoco hace nada contra perturbaciones que varían en el tiempo más rápido que el lazo.
\correcto{El integrador corrige errores constantes; la estabilidad del nuevo equilibrio hay que verificarla aparte (con la linealización en ese equilibrio, o con una función de Lyapunov).}
\end{trampa}

%% ---------------------------------------------------------------------
\section{La idea: cancelar}
\label{sec:fl-idea}
%% ---------------------------------------------------------------------

\subsection{Un péndulo que se vuelve lineal}

Considere el péndulo con un par de entrada,
\begin{equation}
  \dot x_1=x_2,\qquad \dot x_2=-a\sin x_1-b\,x_2+c\,u,
  \label{eq:fl-pendulo}
\end{equation}
con $a=g/\ell$, $b=k/m$ y $c=1/(m\ell^2)$, y la meta de llevarlo a la posición invertida $x_1=\pi$. La no linealidad $-a\sin x_1$ entra en la misma ecuación que la entrada. Entonces se puede \emph{cancelar}:
\begin{equation}
  u=\frac ac\sin x_1+\frac1c\,v
  \quad\Longrightarrow\quad
  \dot x_1=x_2,\quad \dot x_2=-b\,x_2+v .
  \label{eq:fl-cancelapendulo}
\end{equation}
El sistema de la nueva entrada $v$ al estado \emph{es} lineal, en todo el plano, no solo cerca de un punto. Con $v=-k_1(x_1-\pi)-k_2x_2$, el error $e=x_1-\pi$ cumple exactamente $\ddot e+(b+k_2)\dot e+k_1e=0$.

Comparemos con el método 1. Para que la jacobiana en $(\pi,0)$ dé el mismo polinomio característico, $c\,u=-(a+k_1)(x_1-\pi)-k_2x_2$ (porque $-a\sin x_1\approx a(x_1-\pi)$ cerca de $\pi$). Con $a=10$, $b=1$, $c=10$, $k_1=4$, $k_2=3$ (polos en $-2$ y $-2$) y la condición inicial colgante $(0,0)$, la ley exacta reproduce la respuesta lineal diseñada con una diferencia de $10^{-11}$ (el error del integrador numérico); la jacobiana se separa hasta $1.42$\,rad, porque la condición inicial está lejos de donde la aproximación vale (\cref{fig:fl-pendulo}).

\begin{figure}[ht]
\centering
\includegraphics[width=\linewidth]{cap12/f_pendulo}
\caption{Péndulo desde la posición colgante. Izquierda y centro: hacia $\pi$, la ley exacta coincide con el diseño lineal hasta la precisión numérica; la jacobiana no. Derecha: hacia $\pi/2$ con un modelo equivocado ($\hat a=8$ en vez de $a=10$) queda un error de $25.8\grados$.}
\label{fig:fl-pendulo}
\end{figure}

\codigorepo{cap12_linealizacion/C12_Ej3_PenduloExacto.m}{disenos}{Exacta contra jacobiana}

\begin{trampa}[label=tr:fl-robusta]{``Exacta quiere decir perfecta''}
\plausible{La cancelación \eqref{eq:fl-cancelapendulo} hace el sistema lineal exactamente; luego el diseño es inmune a la no linealidad.}
La cancelación es exacta solo si el modelo lo es. Si el controlador usa $\hat a$ en vez de $a$, queda el término $(\hat a-a)\sin x_1$ sin cancelar. Hacia $x_1=\pi/2$ con $\hat a=8$ y $a=10$, el equilibrio del lazo cumple $-2\sin x_1=4(x_1-\pi/2)$, es decir $x_1=1.1206$\,rad: un error de $-25.8\grados$ que no desaparece (\cref{fig:fl-pendulo}, derecha).
\correcto{Linealización exacta no es lo mismo que robustez. Hace falta analizar el error de modelo (y casi siempre agregar acción integral, \cref{sec:fl-levitador}).}
\end{trampa}

\codigorepo{cap12_linealizacion/C12_Ej3_PenduloExacto.m}{robustez}{El mismo diseño con un modelo equivocado}

\begin{pausa}
Con el mismo error de modelo ($\hat a=8$), ¿por qué el péndulo llevado a $\pi$ \emph{no} queda con error permanente?
\respuesta{Porque el término sin cancelar es $(\hat a-a)\sin x_1$, y $\sin\pi=0$: en el equilibrio deseado el error de modelo no produce fuerza. Afecta el transitorio, no el equilibrio. Hacia $\pi/2$, $\sin x_1\approx1$ y el término pesa todo lo que puede.}
\end{pausa}

\subsection{La clase de sistemas y los dos problemas}

¿Cuándo se puede hacer algo parecido? El péndulo tenía una estructura especial: la no linealidad y la entrada estaban en la misma ecuación. En general, la herramienta funciona para sistemas \clave{afines en la entrada}, de una entrada y una salida,
\begin{equation}
  \dot x=f(x)+g(x)\,u,\qquad y=h(x),
  \label{eq:fl-afin}
\end{equation}
con $x\in\R^n$, $u,y\in\R$, y $f,g:D\to\R^n$, $h:D\to\R$ \emph{suaves} (con derivadas parciales continuas de todo orden) en un dominio $D\subset\R^n$. El levitador, el péndulo y van der Pol están en esta clase; casi todos los sistemas mecánicos y electromecánicos con actuador de fuerza o de voltaje también.

\begin{definicion}[label=def:fl-problemas]{Linealización por retroalimentación}
Se busca una ley suave $u=\alpha(x)+\beta(x)\,v$, con $\beta(x)\neq0$, tal que el lazo cerrado
\[
  \dot x=f(x)+g(x)\alpha(x)+g(x)\beta(x)\,v,\qquad y=h(x),
\]
tenga una propiedad lineal. Hay dos problemas distintos:
\begin{enumerate}
  \item \clave{Linealización entrada--salida}: que el mapa de $v$ a $y$ sea lineal (puede quedar una dinámica interna, no lineal, que no se ve en $y$).
  \item \clave{Linealización entrada--estado}: que exista un cambio de coordenadas $z=\Phi(x)$ en el que el lazo cerrado sea $\dot z=Az+Bv$ completo.
\end{enumerate}
\end{definicion}

El péndulo \eqref{eq:fl-cancelapendulo} es linealizable entrada--estado sin cambiar de coordenadas. En general harán falta coordenadas nuevas, y para construirlas, tres herramientas de geometría diferencial.

%% ---------------------------------------------------------------------
\section{Herramientas: derivadas de Lie y difeomorfismos}
\label{sec:fl-herramientas}
%% ---------------------------------------------------------------------

\subsection{Derivadas de Lie}

Un \clave{campo vectorial} suave $f:D\to\R^n$ asigna a cada punto $x$ un vector columna $f(x)=(f_1(x),\dots,f_n(x))$. Un \clave{campo de covectores} asigna a cada punto un vector \emph{fila} $\omega(x)=[\omega_1(x),\dots,\omega_n(x)]$. El ejemplo principal es la \clave{diferencial} de una función suave $\lambda:D\to\R$: el covector $d\lambda(x)=\jac\lambda x(x)$.

\begin{definicion}[label=def:fl-lie]{Derivada de Lie}
La derivada de Lie de $\lambda$ a lo largo de $f$ es la función
\begin{equation}
  L_f\lambda(x)=\jac\lambda x(x)\,f(x)=\sum_{i=1}^n\frac{\partial\lambda}{\partial x_i}(x)\,f_i(x)=\inner{d\lambda(x)}{f(x)} .
  \label{eq:fl-lie}
\end{equation}
Las derivadas sucesivas y mixtas se definen recursivamente:
\[
  L_f^0\lambda=\lambda,\qquad L_f^k\lambda=L_f\bigl(L_f^{k-1}\lambda\bigr),\qquad L_gL_f\lambda=\jac{(L_f\lambda)}x\,g .
\]
\end{definicion}

Geométricamente, $L_f\lambda$ es la proyección de $f$ sobre el gradiente de $\lambda$; si $f$ es tangente a las curvas de nivel de $\lambda$, $L_f\lambda=0$. Y a lo largo de una solución de $\dot x=f(x)$, $\tfrac{d}{dt}\lambda(x(t))=L_f\lambda(x(t))$: la derivada de Lie es la derivada en el tiempo \emph{sin resolver la ecuación}.

\begin{ideaclave}
Ya la usábamos sin nombrarla: para $\dot x=f(x)$, $\Vdot=L_fV$. Toda la teoría de Lyapunov es una colección de desigualdades sobre derivadas de Lie.
\end{ideaclave}

\begin{ejemplo}[label=ej:fl-lievdp]{Van der Pol y su energía}
Para $\dot x_1=x_2$, $\dot x_2=-x_1+\mu(1-x_1^2)x_2$ y $\lambda=\tfrac12(x_1^2+x_2^2)$, calcule $L_f\lambda$ e interprete su signo.
\solucion $L_f\lambda=x_1x_2+x_2\bigl(-x_1+\mu(1-x_1^2)x_2\bigr)=\mu(1-x_1^2)\,x_2^2$. Con $\mu>0$, es no negativa en la franja $\abs{x_1}<1$ y no positiva fuera: el oscilador inyecta energía cerca del origen y la disipa lejos. De ese equilibrio nace el ciclo límite del \cref{cap:orbitas} (\cref{fig:fl-lie}).
\end{ejemplo}

\begin{figure}[ht]
\centering
\includegraphics[width=\linewidth]{cap12/f_lie}
\caption{Van der Pol con $\mu=1$. Izquierda: la franja $\abs{x_1}<1$ donde $L_f\lambda\ge0$, círculos de nivel de $\lambda$ y una solución que llega al ciclo. Derecha: $L_f\lambda$ calculada con la fórmula coincide con $d\lambda/dt$ medida sobre la solución.}
\label{fig:fl-lie}
\end{figure}

\codigorepo{cap12_linealizacion/C12_Ej4_DerivadasLie.m}{vdp}{Derivada de Lie: fórmula contra derivada medida}

\subsection{Difeomorfismos}

\begin{definicion}[label=def:fl-difeo]{Difeomorfismo}
Una función $\Phi:D\to\R^n$, con $D\subset\R^n$ abierto, es un \clave{difeomorfismo} de $D$ en $\Phi(D)$ si es uno a uno, suave, y su inversa $\Phi^{-1}:\Phi(D)\to D$ también es suave. Es un \clave{difeomorfismo local} alrededor de $x^o$ si lo es en alguna vecindad de $x^o$, y un \clave{difeomorfismo global} si $D=\Phi(D)=\R^n$.
\end{definicion}

Un difeomorfismo es un cambio de coordenadas no lineal, como lo es una matriz invertible $T$ en $z=Tx$ para un sistema lineal. Reconocer uno local es fácil.

\begin{teorema}[label=thm:fl-inversa]{Función inversa}
Si $\Phi$ es suave en $D$ y el jacobiano $\jac\Phi x$ es no singular en $x^o\in D$, entonces $\Phi$ es un difeomorfismo local alrededor de $x^o$. Como las filas de $\jac\Phi x$ son los covectores $d\phi_1(x),\dots,d\phi_n(x)$, la condición equivale a que esos covectores sean linealmente independientes en $x^o$ (y entonces, por continuidad, en una vecindad).
\end{teorema}

\begin{trampa}[label=tr:fl-difeo]{``Jacobiano no singular en todo punto, difeomorfismo global''}
\plausible{Si $\det\jac\Phi x\neq0$ en todo $\R^n$, el \cref{thm:fl-inversa} se aplica en todas partes, y $\Phi$ es un difeomorfismo global.}
El teorema da un difeomorfismo \emph{alrededor de cada punto}, no uno en todo el espacio. Para $\Psi(x)=e^{x_1}(\cos x_2,\ \sin x_2)$, $\det\jac\Psi x=e^{2x_1}>0$ en todo $\R^2$, pero $\Psi(x_1,x_2+2\pi)=\Psi(x_1,x_2)$: no es uno a uno. La franja $\{0\le x_2<4\pi\}$ cubre el anillo dos veces (\cref{fig:fl-difeo}, derecha).
\correcto{Para un difeomorfismo global hace falta algo más. El criterio de Hadamard: si $\Phi:\R^n\to\R^n$ es suave, su jacobiano es no singular en todo punto \emph{y} $\norm{\Phi(x)}\to\infty$ cuando $\norm x\to\infty$, entonces $\Phi$ es un difeomorfismo global. $\Psi$ falla la segunda condición: con $x_1\to-\infty$, $\Psi\to0$.}
\end{trampa}

\begin{figure}[ht]
\centering
\includegraphics[width=\linewidth]{cap12/f_difeo}
\caption{Izquierda: la imagen de una rejilla bajo el difeomorfismo global $\Phi(x)=(x_1,\ x_2+x_1^3/2)$. Derecha: bajo $\Psi(x)=e^{x_1}(\cos x_2,\sin x_2)$, $x_2\in[0,2\pi)$ (continua) y $x_2\in[2\pi,4\pi)$ (discontinua) dan las mismas circunferencias.}
\label{fig:fl-difeo}
\end{figure}

\video{V12_4_difeomorfismo}{Local en todo punto no es global}

\codigorepo{cap12_linealizacion/C12_Ej5_Difeomorfismos.m}{local}{Jacobiano no singular, función no inyectiva}

\subsection{Cambio de coordenadas}

Sea $z=\Phi(x)$ un difeomorfismo en $D$. Por la regla de la cadena, $\dot z=\jac\Phi x\,\dot x$, y el sistema afín \eqref{eq:fl-afin} sigue siendo afín en las nuevas coordenadas:

\begin{proposicion}[label=prop:fl-coordenadas]{El sistema en coordenadas nuevas}
En $z=\Phi(x)=(\phi_1(x),\dots,\phi_n(x))$, el sistema \eqref{eq:fl-afin} se escribe $\dot z=\bar f(z)+\bar g(z)\,u$, $y=\bar h(z)$, en $\Phi(D)$, con
\begin{equation}
  \bar f(z)=\begin{bmatrix}L_f\phi_1(x)\\ \vdots\\ L_f\phi_n(x)\end{bmatrix}_{x=\Phi^{-1}(z)},\qquad
  \bar g(z)=\begin{bmatrix}L_g\phi_1(x)\\ \vdots\\ L_g\phi_n(x)\end{bmatrix}_{x=\Phi^{-1}(z)},\qquad
  \bar h(z)=h\bigl(\Phi^{-1}(z)\bigr) .
  \label{eq:fl-coordenadas}
\end{equation}
\end{proposicion}

Cada coordenada nueva evoluciona según su derivada de Lie: $\dot z_i=\dot\phi_i=L_f\phi_i+L_g\phi_i\,u$. El truco de todo lo que sigue es escoger las $\phi_i$ para que esas derivadas tengan una forma sencilla.

\begin{ejemplo}[label=ej:fl-coordvdp]{Van der Pol en otras coordenadas}
Escriba van der Pol (con entrada, $\dot x_2=-x_1+\mu(1-x_1^2)x_2+u$) en las coordenadas $z=(x_1^2+x_2,\ x_1)$.
\solucion El jacobiano $\begin{bmatrix}2x_1&1\\1&0\end{bmatrix}$ tiene determinante $-1$ en todo punto, y la inversa es explícita: $x=\Phi^{-1}(z)=(z_2,\ z_1-z_2^2)$, así que $\Phi$ es un difeomorfismo global. Con \eqref{eq:fl-coordenadas}: $L_g\phi_1=1$, $L_g\phi_2=0$, $L_f\phi_2=x_2=z_1-z_2^2$ y $L_f\phi_1=2x_1x_2-x_1+\mu(1-x_1^2)x_2$. Luego
\[
  \dot z_1=2z_2(z_1-z_2^2)-z_2+\mu(1-z_2^2)(z_1-z_2^2)+u,\qquad \dot z_2=z_1-z_2^2 .
\]
La entrada aparece solo en $\dot z_1$. En la \cref{sec:fl-grado} veremos que no es casualidad.
\end{ejemplo}

\codigorepo{cap12_linealizacion/C12_Ej5_Difeomorfismos.m}{vdp}{La inversa del cambio de coordenadas}

%% ---------------------------------------------------------------------
\section{Grado relativo}
\label{sec:fl-grado}
%% ---------------------------------------------------------------------

\subsection{Derivar la salida hasta que aparezca la entrada}

Derivemos la salida de \eqref{eq:fl-afin} a lo largo de las soluciones:
\[
  \dot y=\jac hx\bigl(f(x)+g(x)u\bigr)=L_fh(x)+L_gh(x)\,u .
\]
Si $L_gh\equiv0$ cerca del punto de interés, la entrada no aparece, $\dot y=L_fh(x)$, y podemos derivar otra vez:
\[
  \ddot y=L_f^2h(x)+L_gL_fh(x)\,u .
\]
Si también $L_gL_fh\equiv0$, seguimos. Si en la derivada $r$-ésima aparece por fin la entrada,
\begin{equation}
  y^{(r)}=L_f^rh(x)+L_gL_f^{r-1}h(x)\,u,\qquad L_gL_f^{r-1}h(x)\neq0 .
  \label{eq:fl-yr}
\end{equation}

\begin{definicion}[label=def:fl-grado]{Grado relativo \citep[\S4.1]{Isi95}}
El sistema \eqref{eq:fl-afin} tiene \clave{grado relativo} $r$ en $x^o$ si
\begin{enumerate}
  \item $L_gL_f^ih(x)=0$ para todo $x$ en una vecindad de $x^o$ y todo $i=0,1,\dots,r-2$, y
  \item $L_gL_f^{r-1}h(x^o)\neq0$.
\end{enumerate}
\end{definicion}

\begin{ideaclave}
El grado relativo es el número de veces que hay que derivar la salida para que la entrada aparezca explícitamente. Entre $u$ y $y$ hay $r$ ``integradores'': la entrada afecta directamente a $y^{(r)}$, y a $y$ solo después de $r$ integraciones.
\end{ideaclave}

Observe que la condición 1 se pide en toda una \emph{vecindad}, y la 2 solo en el punto (por continuidad, vale entonces en una vecindad). El grado relativo depende de la salida $h$: es una propiedad del par planta--salida, no de la planta sola. Y no puede pasar de $n$ (\cref{lem:fl-independencia}).

\begin{trampa}[label=tr:fl-nodefinido]{``Todo sistema tiene grado relativo en todo punto''}
\plausible{Basta derivar $y$ hasta que aparezca $u$; en algún momento aparece, y ese es $r$.}
El grado relativo puede \emph{no existir} en un punto. Ocurre cuando la primera función de la sucesión $L_gh,\ L_gL_fh,\ L_gL_f^2h,\dots$ que no es idénticamente nula cerca de $x^o$ se anula justo en $x^o$. El ejemplo mínimo: $\dot x=u$, $y=x^2$. Como $L_gh=2x$, el grado relativo es $1$ en todo $x\neq0$; en $x^o=0$, $L_gh(0)=0$ (falla la condición 2 con $r=1$) pero $L_gh$ no es idénticamente nula cerca de $0$ (falla la condición 1 con $r\ge2$). No hay grado relativo en el origen.
\correcto{El conjunto de puntos donde hay grado relativo es abierto y denso~\citep{Isi95}, pero puede excluir justo el punto que interesa. En el levitador excluye $i=0$ (\cref{ej:fl-gradolev}), y eso tiene consecuencias físicas.}
\end{trampa}

\subsection{El caso lineal}

\begin{proposicion}[label=prop:fl-gradolineal]{Grado relativo de un sistema lineal}
Para $\dot x=Ax+Bu$, $y=Cx$ (es decir, $f=Ax$, $g=B$, $h=Cx$), $L_f^kh(x)=CA^kx$ y $L_gL_f^kh=CA^kB$. El grado relativo está definido en \emph{todo} punto y es el entero $r$ con
\[
  CA^kB=0\ \ (0\le k\le r-2),\qquad CA^{r-1}B\neq0 .
\]
Coincide con la diferencia entre el grado del denominador y el del numerador de $G(s)=C(sI-A)^{-1}B$.
\end{proposicion}

\begin{proof}
Las fórmulas de las derivadas de Lie se obtienen por inducción: $L_f(CA^kx)=CA^k\,Ax$. Como $L_gL_f^kh$ es constante, las condiciones de la \cref{def:fl-grado} no dependen del punto. Para la última afirmación, $G(s)=\sum_{k\ge1}CA^{k-1}B\,s^{-k}$ para $\abs s$ grande, de modo que $G(s)=CA^{r-1}B\,s^{-r}+O(s^{-r-1})$: el primer término no nulo dice cuánto más grande es el grado del denominador.
\end{proof}

Los números $CA^{k-1}B$ son los \clave{parámetros de Markov}; no cambian con un cambio de base $z=Tx$. Hay una lectura en el tiempo: la respuesta al escalón es $y(t)=\sum_{k\ge1}CA^{k-1}B\,t^k/k!$, así que cerca de $t=0$, $y(t)\approx CA^{r-1}B\,t^r/r!$, y en escala log-log la pendiente es $r$ (\cref{fig:fl-gradolineal}).

\begin{figure}[ht]
\centering
\includegraphics[width=0.66\linewidth]{cap12/f_grado_lineal}
\caption{Respuestas al escalón cerca de $t=0$ en escala log-log: la pendiente es el grado relativo. La punteada es $t^r/r!$. El cero de $(s+2)/((s+1)(s+3))$ no cambia la pendiente: $r=1$.}
\label{fig:fl-gradolineal}
\end{figure}

\codigorepo{cap12_linealizacion/C12_Ej6_GradoRelativo.m}{markov}{Parámetros de Markov y pendiente de la respuesta al escalón}

\subsection{Ejemplos no lineales}

\begin{ejemplo}[label=ej:fl-gradovdp]{Van der Pol con dos salidas}
Calcule el grado relativo de $\dot x_1=x_2$, $\dot x_2=-x_1+\mu(1-x_1^2)x_2+u$ con (a) $y=x_1$ y (b) $y=x_1^2+x_2$.
\solucion (a) $L_gh=\jac{x_1}x\,g=[1\ \ 0]\,[0\ \ 1]\T=0$, $L_fh=x_2$ y $L_gL_fh=1$: $r=2=n$ en todo $\R^2$. (b) $L_gh=[2x_1\ \ 1]\,[0\ \ 1]\T=1$: $r=1<n$ en todo $\R^2$.
\end{ejemplo}

\begin{trampa}[label=tr:fl-salida]{``El grado relativo es del sistema''}
\plausible{Van der Pol tiene grado relativo $2$.}
La misma planta tiene $r=2$ con $y=x_1$ y $r=1$ con $y=x_1^2+x_2$. Lo mismo pasa con la dinámica cero (\cref{sec:fl-cero}).
\correcto{El grado relativo es una propiedad del par (planta, salida). Cambiar lo que se mide, o lo que se quiere regular, cambia todo el análisis.}
\end{trampa}

\begin{ejemplo}[label=ej:fl-gradolev]{El levitador}
Calcule el grado relativo del levitador \eqref{eq:fl-levitador} con $y=x_1$.
\solucion Aquí $g(x)=\bigl(0,\,0,\,1/L(x_1)\bigr)$. Primero $L_gh=0$ y $L_fh=x_2$; luego $L_gL_fh=0$ y
\[
  L_f^2h=g-\frac km\,x_2-\frac{L_0a\,x_3^2}{2m(a+x_1)^2} ,
\]
la aceleración. Como esta depende de $x_3$,
\begin{equation}
  L_gL_f^2h(x)=-\frac{L_0\,a\,x_3}{m\,(a+x_1)^2\,L(x_1)} .
  \label{eq:fl-blev}
\end{equation}
El grado relativo es $r=3=n$ en todo punto con $x_3\neq0$; en el equilibrio de $\bar y=0.05$, $L_gL_f^2h=-125.28$. En $x_3=0$ no está definido (\cref{fig:fl-inestable}, derecha).
\end{ejemplo}

\begin{pausa}
Explique físicamente por qué el levitador pierde el grado relativo justo cuando la corriente es cero.
\respuesta{La fuerza es $F\propto-i^2$, así que $\partial F/\partial i=0$ en $i=0$: un cambio pequeño de corriente, que es lo que el voltaje cambia primero, no cambia la fuerza a primer orden. La entrada pierde su efecto sobre la aceleración. Detrás está algo más elemental: el electroimán solo atrae, y con $i=0$ ni siquiera eso.}
\end{pausa}

\codigorepo{cap12_linealizacion/C12_Ej6_GradoRelativo.m}{levitador}{El grado relativo del levitador y su singularidad}

\subsection{Por qué $r\le n$}

\begin{lema}[label=lem:fl-independencia]{Independencia \citep[Lema~4.1.1]{Isi95}}
Si el sistema tiene grado relativo $r$ en $x^o$, los covectores
\[
  dh(x),\ dL_fh(x),\ \dots,\ dL_f^{r-1}h(x)
\]
son linealmente independientes para todo $x$ en una vecindad de $x^o$. En particular, $r\le n$.
\end{lema}

\begin{proof}[Idea de la prueba]
Con el corchete de Lie $[f,g]=\jac gxf-\jac fxg$ y sus iterados $\mathrm{ad}_f^jg=[f,\mathrm{ad}_f^{j-1}g]$, se prueba por inducción la identidad
\[
  L_{\mathrm{ad}_f^jg}L_f^ih(x)=\begin{cases}0,& i+j\le r-2,\\ (-1)^j\,L_gL_f^{r-1}h(x),& i+j=r-1,\end{cases}
\]
en una vecindad de $x^o$. Entonces la matriz $r\times r$ de entradas $\inner{dL_f^ih(x)}{\mathrm{ad}_f^jg(x)}$ es triangular respecto a la antidiagonal, con antidiagonal $\pm L_gL_f^{r-1}h\neq0$: es no singular. Si una combinación $\sum_ic_i\,dL_f^ih$ fuera cero, al aparearla con $g,\mathrm{ad}_fg,\dots$ daría un sistema lineal con esa matriz, y $c=0$.
\end{proof}

Por el \cref{thm:fl-inversa}, las $r$ funciones $h,L_fh,\dots,L_f^{r-1}h$ son candidatas a ser parte de un sistema de coordenadas. Faltan $n-r$: el análogo de completar una base en álgebra lineal.

%% ---------------------------------------------------------------------
\section{Coordenadas adaptadas a la salida y forma normal}
\label{sec:fl-forma}
%% ---------------------------------------------------------------------

\begin{teorema}[label=thm:fl-coordenadas]{Coordenadas adaptadas \citep[Prop.~4.1.3]{Isi95}}
Suponga que el sistema \eqref{eq:fl-afin} tiene grado relativo $r$ en $x^o$, y sean
\[
  \phi_1(x)=h(x),\quad \phi_2(x)=L_fh(x),\quad\dots,\quad \phi_r(x)=L_f^{r-1}h(x) .
\]
Si $r<n$, existen $n-r$ funciones suaves $\phi_{r+1},\dots,\phi_n$ tales que $\Phi=(\phi_1,\dots,\phi_n)$ tiene jacobiano no singular en $x^o$ (un difeomorfismo local alrededor de $x^o$). Además, siempre se pueden escoger de modo que
\begin{equation}
  L_g\phi_i(x)=0,\qquad i=r+1,\dots,n,
  \label{eq:fl-Lgphi}
\end{equation}
en una vecindad de $x^o$.
\end{teorema}

La primera parte es completar una base, gracias al \cref{lem:fl-independencia}. La segunda dice algo más fino: las coordenadas que faltan se pueden escoger para que $u$ no las mueva. Encontrarlas exige resolver las ecuaciones en derivadas parciales $L_g\phi_i=\jac{\phi_i}xg=0$; en los ejemplos se adivinan, y no son únicas.

\subsection{La forma normal}

En las coordenadas del \cref{thm:fl-coordenadas}, separemos $z=\Phi(x)=(\xi,\eta)$:
\begin{equation}
  \xi=\begin{bmatrix}h(x)\\ L_fh(x)\\ \vdots\\ L_f^{r-1}h(x)\end{bmatrix}\in\R^r,
  \qquad
  \eta=\begin{bmatrix}\phi_{r+1}(x)\\ \vdots\\ \phi_n(x)\end{bmatrix}\in\R^{n-r}.
  \label{eq:fl-xieta}
\end{equation}
Por la \cref{prop:fl-coordenadas}, $\dot\xi_i=L_f\phi_i+L_g\phi_i\,u$. Para $i<r$, $L_g\phi_i=L_gL_f^{i-1}h=0$ y $L_f\phi_i=L_f^ih=\phi_{i+1}$; para $i=r$ aparece la entrada; para las $\eta$, \eqref{eq:fl-Lgphi} elimina la entrada. El resultado es la \clave{forma normal}:
\begin{equation}
  \begin{aligned}
    \dot\xi_1&=\xi_2,\quad\dots,\quad\dot\xi_{r-1}=\xi_r,\\
    \dot\xi_r&=a(\xi,\eta)+b(\xi,\eta)\,u,\\
    \dot\eta&=q(\xi,\eta),\\
    y&=\xi_1,
  \end{aligned}
  \qquad\text{con}\qquad
  \begin{aligned}
    a(z)&=L_f^rh(x)\big|_{x=\Phi^{-1}(z)},\\
    b(z)&=L_gL_f^{r-1}h(x)\big|_{x=\Phi^{-1}(z)},\\
    q_i(z)&=L_f\phi_i(x)\big|_{x=\Phi^{-1}(z)} .
  \end{aligned}
  \label{eq:fl-formanormal}
\end{equation}
En forma compacta (forma de Brunovsky para la parte $\xi$),
\begin{equation}
  \dot\xi=A_c\xi+B_c\bigl[a(\xi,\eta)+b(\xi,\eta)u\bigr],\qquad y=C_c\xi,
  \label{eq:fl-brunovsky}
\end{equation}
con $A_c$ la matriz $r\times r$ con unos sobre la diagonal y ceros en el resto, $B_c=(0,\dots,0,1)$ y $C_c=[1\ 0\ \cdots\ 0]$. Las primeras $r$ coordenadas \emph{son} la salida y sus derivadas: $\xi_1(t)=y(t)$, $\xi_2(t)=\dot y(t)$, \dots, $\xi_r(t)=y^{(r-1)}(t)$.

\begin{figure}[ht]
\centering
\includegraphics[width=0.85\linewidth]{cap12/t_forma_normal}
\caption{La forma normal: una cadena de $r$ integradores, con la entrada solo en la primera, y una dinámica interna $\eta$ que recibe $\xi$ pero no la entrada.}
\label{fig:fl-formanormal}
\end{figure}

\notamargen{\textbf{Khalil y estas notas.} Khalil escribe $\rho$ para el grado relativo, $\dot\eta=f_0(\eta,\xi)$ para la dinámica interna y $\gamma(x)\bigl[u-\alpha(x)\bigr]$ para $a+bu$. Aquí seguimos a Isidori: $r$, $q$, $a$ y $b$.}

\begin{ejemplo}[label=ej:fl-formaexp]{Una forma normal con $r=2<n=3$ \citep{Isi95}}
Lleve a la forma normal
\[
  \dot x=\begin{bmatrix}-x_1\\ x_1x_2\\ x_2\end{bmatrix}+\begin{bmatrix}e^{x_2}\\ 1\\ 0\end{bmatrix}u,\qquad y=x_3 .
\]
\solucion $L_gh=0$, $L_fh=x_2$ y $L_gL_fh=1$: $r=2$ en todo $\R^3$. Tomamos $\xi_1=x_3$, $\xi_2=x_2$ y buscamos $\phi_3$ con $L_g\phi_3=e^{x_2}\partial_{x_1}\phi_3+\partial_{x_2}\phi_3=0$. Sirve $\phi_3=1+x_1-e^{x_2}$. El jacobiano de $\Phi=(x_3,\,x_2,\,1+x_1-e^{x_2})$ tiene determinante $-1$ y la inversa es explícita ($x_1=-1+\eta+e^{\xi_2}$, $x_2=\xi_2$, $x_3=\xi_1$): es un difeomorfismo global. Con $L_f^2h=x_1x_2$ y $L_f\phi_3=-x_1-e^{x_2}x_1x_2$,
\[
  \dot\xi_1=\xi_2,\qquad
  \dot\xi_2=\bigl(-1+\eta+e^{\xi_2}\bigr)\xi_2+u,\qquad
  \dot\eta=-\bigl(-1+\eta+e^{\xi_2}\bigr)\bigl(1+\xi_2e^{\xi_2}\bigr) .
\]
Note el signo de la última ecuación (la versión anterior de estas notas lo había perdido): con $\xi=0$ queda $\dot\eta=-\eta$.
\end{ejemplo}

\begin{ejemplo}[label=ej:fl-formakhalil]{Otra, con un difeomorfismo global}
Para
\begin{equation}
  \dot x=\begin{bmatrix}x_3-x_2^3\\ -x_2\\ x_1^2-x_3\end{bmatrix}+\begin{bmatrix}0\\-1\\1\end{bmatrix}u,\qquad y=x_1,
  \label{eq:fl-ejkhalil}
\end{equation}
$L_gh=0$, $L_fh=x_3-x_2^3$ y $L_gL_fh=1+3x_2^2\neq0$: $r=2$ en todo $\R^3$. Con $\phi_3=x_2+x_3$ ($L_g\phi_3=-1+1=0$), $\Phi=(x_1,\ x_3-x_2^3,\ x_2+x_3)$ tiene $\det\jac\Phi x=-(1+3x_2^2)\neq0$, y es global: dado $z$, $x_2$ es la única raíz real de $x_2+x_2^3=z_3-z_2$. La forma normal es
\[
  \dot z_1=z_2,\qquad \dot z_2=a(z)+b(z)u,\qquad \dot z_3=z_1^2-z_3,
\]
con $a=L_f^2h=3x_2^3+x_1^2-x_3$ y $b=1+3x_2^2$ evaluadas en $x=\Phi^{-1}(z)$. Volveremos a este ejemplo en la \cref{sec:fl-estabilidad}.
\end{ejemplo}

\codigorepo{cap12_linealizacion/C12_Ej7_FormaNormal.m}{isidori}{La forma normal y el sistema original dan la misma solución}

\begin{cuidado}[Las $\phi_i$ complementarias no son únicas]
Cualquier función de $\phi_3$ que conserve el difeomorfismo sirve igual en el \cref{ej:fl-formaexp}. Lo que \emph{sí} es único es la dinámica que queda cuando la salida se anula (\cref{sec:fl-cero}): no depende de cuáles $\phi_i$ se escojan.
\end{cuidado}

%% ---------------------------------------------------------------------
\section{Linealización entrada--salida}
\label{sec:fl-io}
%% ---------------------------------------------------------------------

\subsection{La ley que linealiza}

Si el sistema tiene grado relativo $r$ en $x^o$, $b(z^o)=L_gL_f^{r-1}h(x^o)\neq0$ y, por continuidad, $b\neq0$ en una vecindad. Ahí se puede despejar la entrada de \eqref{eq:fl-formanormal}:
\begin{equation}
  u=\frac{1}{b(z)}\bigl[-a(z)+v\bigr]
  \qquad\text{es decir}\qquad
  u=\frac{1}{L_gL_f^{r-1}h(x)}\Bigl[-L_f^rh(x)+v\Bigr] .
  \label{eq:fl-ley}
\end{equation}
Es de la forma $u=\alpha(x)+\beta(x)v$, con $\alpha=-L_f^rh/L_gL_f^{r-1}h$ y $\beta=1/L_gL_f^{r-1}h$, y no hace falta calcular $\Phi^{-1}$ para aplicarla: basta evaluar derivadas de Lie en el estado medido. El lazo cerrado es
\begin{equation}
  \dot\xi=A_c\xi+B_cv,\qquad y=C_c\xi,\qquad \dot\eta=q(\xi,\eta),
  \label{eq:fl-lazocerrado}
\end{equation}
y el mapa de $v$ a $y$ es una cadena de $r$ integradores: $G(s)=C_c(sI-A_c)^{-1}B_c=1/s^r$.

\begin{ideaclave}
La realimentación volvió $\eta$ \emph{no observable} desde $y$: la salida obedece $y^{(r)}=v$ sin importar lo que haga $\eta$. Que $y$ se comporte perfectamente no dice nada de $\eta$.
\end{ideaclave}

\subsection{El lazo externo}

Sobre la cadena de integradores se diseña un controlador lineal.

\begin{proposicion}[label=prop:fl-seguimiento]{Estabilización y seguimiento}
Sean $k_1,\dots,k_r$ tales que $s^r+k_rs^{r-1}+\dots+k_2s+k_1$ es Hurwitz, y $K=[k_1\ \cdots\ k_r]$.
\begin{enumerate}
  \item (Estabilización.) Con $v=-K\xi$, $\dot\xi=(A_c-B_cK)\xi$ y $\xi\to0$ exponencialmente.
  \item (Seguimiento.) Sea $y_d(t)$ una referencia con $r$ derivadas continuas, y $\bar y_d=(y_d,\dot y_d,\dots,y_d^{(r-1)})$. Con
  \begin{equation}
    v=y_d^{(r)}-K\bigl(\xi-\bar y_d\bigr),
    \label{eq:fl-vseguimiento}
  \end{equation}
  el error $e=\xi-\bar y_d$ cumple $\dot e=(A_c-B_cK)e$, y $y(t)-y_d(t)\to0$ exponencialmente.
\end{enumerate}
\end{proposicion}

\begin{proof}
En 2, $\dot e_i=\xi_{i+1}-y_d^{(i)}=e_{i+1}$ para $i<r$, y $\dot e_r=v-y_d^{(r)}=-Ke$.
\end{proof}

La \cref{prop:fl-seguimiento} habla solo de $\xi$. Si además $\dot\eta=q(\xi,\eta)$ es ISS con $\xi$ como entrada y las derivadas de $y_d$ están acotadas, entonces $\eta$ también queda acotada y \emph{todas} las señales internas lo están~\citep{Arc18}. Sin esa hipótesis, la salida puede seguir la referencia mientras el estado interno diverge (\cref{sec:fl-cero}).

\begin{figure}[ht]
\centering
\includegraphics[width=0.95\linewidth]{cap12/t_lazo_fl}
\caption{Dos lazos. El interno (no lineal) cancela con $u=(v-a(z))/b(z)$ y deja de $v$ a $y$ una cadena de $r$ integradores; el externo (lineal) estabiliza o hace seguir la referencia.}
\label{fig:fl-lazo}
\end{figure}

\begin{trampa}[label=tr:fl-escalon]{``Para seguir un escalón basta poner $y_d$ constante a trozos''}
\plausible{La ley \eqref{eq:fl-vseguimiento} sigue cualquier referencia; un escalón es la más sencilla.}
La ley pide $y_d,\dot y_d,\dots,y_d^{(r)}$, y un escalón no tiene derivadas en el instante del salto: con $r\ge1$ no es una referencia admisible. En la práctica, conectarlo directamente produce en $v$ (y en $u$) un salto proporcional a $k_r,\dots,k_1$ y al tamaño del escalón, que el actuador puede no tener, y que en el levitador lleva a la singularidad de la \cref{sec:fl-levitador}.
\correcto{Se planifica una trayectoria suave: por ejemplo, se pasa el escalón por un filtro $p^{r+1}/(s+p)^{r+1}$, que entrega $y_d$ y sus derivadas hasta la $r$-ésima.}
\end{trampa}

\subsection{Si $r=n$: linealización entrada--estado}

Si el grado relativo es $n$, no hay $\eta$: $\Phi(x)=(h,L_fh,\dots,L_f^{n-1}h)$ es el cambio de coordenadas completo, y la ley
\[
  u=\frac{1}{L_gL_f^{n-1}h(x)}\bigl[v-L_f^nh(x)\bigr]
\]
da el sistema lineal $\dot z=A_cz+B_cv$ de dimensión $n$. Es la linealización entrada--estado de la \cref{def:fl-problemas}; el levitador con $y=x_1$ es de esta clase.

\begin{ejemplo}[label=ej:fl-vdpfl]{Van der Pol linealizado}
Con $y=x_1$, van der Pol tiene $r=2=n$ (\cref{ej:fl-gradovdp}) y ya está en forma normal, con $a(x)=-x_1+\mu(1-x_1^2)x_2$ y $b=1$. La ley
\[
  u=x_1-\mu(1-x_1^2)x_2+v,\qquad v=-4x_1-4x_2,
\]
da $\ddot y+4\dot y+4y=0$ (polos en $-2$, $-2$). Como $\Phi$ es la identidad, el resultado es \emph{global}: el ciclo límite desaparece por construcción. Con $\mu=1$, el oscilador libre hasta $t=10$ y la ley conectada desde ahí, la salida sigue la predicción lineal con una diferencia menor que $10^{-9}$ (\cref{fig:fl-vdpfl}).
\end{ejemplo}

\begin{figure}[ht]
\centering
\includegraphics[width=\linewidth]{cap12/f_vdp_fl}
\caption{Van der Pol libre (gris) hasta $t=10$; desde ahí, con la ley linealizante. La salida coincide con la solución de $\ddot y+4\dot y+4y=0$.}
\label{fig:fl-vdpfl}
\end{figure}

\codigorepo{cap12_linealizacion/C12_Ej8_VanDerPolFL.m}{ley}{Van der Pol: del ciclo límite a un sistema lineal}

%% ---------------------------------------------------------------------
\section{El levitador, exacto}
\label{sec:fl-levitador}
%% ---------------------------------------------------------------------

Con $y=x_1$ el levitador tiene $r=3=n$ donde $x_3\neq0$ (\cref{ej:fl-gradolev}). Las coordenadas son posición, velocidad y aceleración,
\begin{equation}
  \xi=\Phi(x)=\begin{bmatrix}x_1\\ x_2\\ g-\dfrac km\,x_2-\dfrac{L_0\,a\,x_3^2}{2m(a+x_1)^2}\end{bmatrix},
  \label{eq:fl-philev}
\end{equation}
que es un difeomorfismo en $\{x_3>0\}$ (para $x_1$ y $x_2$ dados, $\xi_3$ es estrictamente decreciente en $x_3>0$). La ley es
\begin{equation}
  u=\frac{v-L_f^3h(x)}{L_gL_f^2h(x)},\qquad
  v=y_d^{(3)}-k_1(\xi_1-y_d)-k_2(\xi_2-\dot y_d)-k_3(\xi_3-\ddot y_d),
  \label{eq:fl-leylev}
\end{equation}
con $L_gL_f^2h$ de \eqref{eq:fl-blev} y, escribiendo $c=L_0a/(2m)$ y $\dot x_3=f_3(x)$ cuando $u=0$,
\[
  L_f^3h=\frac{2c\,x_3^2}{(a+x_1)^3}\,x_2-\frac km\,\xi_3-\frac{2c\,x_3}{(a+x_1)^2}\,f_3(x),
  \qquad
  f_3(x)=\frac{1}{L(x_1)}\Bigl(-Rx_3+\frac{L_0a\,x_2x_3}{(a+x_1)^2}\Bigr).
\]
Usamos $s^3+60s^2+1100s+6000=(s+10)(s+20)(s+30)$, es decir $K=[6000,\ 1100,\ 60]$. No hay dinámica interna: todo el estado es la salida y sus derivadas.

\subsection{Seguimiento: exacta contra jacobiana}

La referencia $y_d=0.08+0.03\sin(2\pi t)$\,m recorre de $5$ a $11$\,cm. La ley jacobiana usa la ganancia fija diseñada en $\bar y=0.05$ con la prealimentación $u^*(y_d(t))$ de \eqref{eq:fl-equilibrio}; la exacta usa \eqref{eq:fl-leylev}. Para $t>1$\,s, la jacobiana tiene un error de $2.36$\,cm, con desfase; la exacta, $2\times10^{-6}$\,m, que es el error numérico (\cref{fig:fl-seguimiento}). La linealización exacta conserva la respuesta diseñada en \emph{todo} el rango de operación; la jacobiana, solo cerca del punto de diseño.

\begin{figure}[ht]
\centering
\includegraphics[width=\linewidth]{cap12/f_seguimiento}
\caption{Seguimiento de $y_d=0.08+0.03\sin 2\pi t$. Izquierda: posiciones. Derecha: error en escala logarítmica; el de la ley exacta decae como $e^{-10t}$ hasta el piso numérico.}
\label{fig:fl-seguimiento}
\end{figure}

\video{V12_1_levitador}{Dos levitadores siguen la misma referencia}

\codigorepo{cap12_linealizacion/C12_Ej9_LevitadorExacto.m}{seguimiento}{Seguimiento con la ley exacta y con la jacobiana}

\subsection{La singularidad}

\begin{trampa}[label=tr:fl-singular]{``La ley exacta lleva la esfera a cualquier posición''}
\plausible{Con \eqref{eq:fl-leylev} y $y_d$ constante, el lazo cerrado es lineal y estable: un escalón de $0.05$ a $0.2$\,m da una respuesta con polos en $-10$, $-20$ y $-30$.}
Para caer $15$\,cm con esos polos, la esfera tendría que acelerar hacia abajo más que $g$, y el electroimán solo atrae: la única forma de acercarse es \emph{apagar} la corriente. La ley lo intenta: $i$ baja a $0.02$\,A en $20$\,ms. Pero $b(x)=L_gL_f^2h\propto x_3$, de modo que $b\to0$ y $u=(v-L_f^3h)/b$ explota (\cref{fig:fl-singularidad}, curvas rojas). La matemática ve lo que la física impone.
\correcto{La ley solo vale donde existe el grado relativo, y la trayectoria tiene que quedarse allí. Se planifica una referencia alcanzable: el escalón pasado por $6^4/(s+6)^4$ da $y_d$ con $\ddot y_d\le0.705$\,m/s$^2<g$; la esfera llega a $0.19966$\,m en $t=2$\,s con una corriente que nunca baja de $6.15$\,A (curvas azules).}
\end{trampa}

\begin{figure}[ht]
\centering
\includegraphics[width=\linewidth]{cap12/f_singularidad}
\caption{Un escalón de $5$ a $20$\,cm (rojo) lleva la corriente a cero en $20$\,ms, donde $b(x)=0$; la misma maniobra con una referencia suave (azul) mantiene la corriente lejos de cero.}
\label{fig:fl-singularidad}
\end{figure}

\codigorepo{cap12_linealizacion/C12_Ej9_LevitadorExacto.m}{suave}{Una referencia suave evita la singularidad}

\subsection{Error de modelo}

Desde $t=0.5$\,s la masa real es $0.12$\,kg y la ley sigue cancelando con $m=0.10$. Sola, la ley exacta queda con un error permanente de $1.96$\,cm, \emph{peor} que el de la realimentación de estados sobre la jacobiana ($1.03$\,cm, \cref{sec:fl-jacobiana}): la cancelación de la gravedad estaba ``perfectamente'' calculada con la masa equivocada. El remedio es el mismo de la \cref{sec:fl-jacobiana}: se agrega $\dot\sigma=y-\bar y$ al comienzo de la cadena y se diseña sobre $(\sigma,\xi_1,\xi_2,\xi_3)$,
\[
  v=-k_0\,\sigma-k_1(\xi_1-\bar y)-k_2\,\xi_2-k_3\,\xi_3,
\]
con $s^4+k_3s^3+k_2s^2+k_1s+k_0=(s+10)(s+15)(s+20)(s+30)$, es decir $(k_3,k_2,k_1,k_0)=(75,\ 2000,\ 22500,\ 90000)$. El error se anula (\cref{fig:fl-robustez}).

\begin{figure}[ht]
\centering
\includegraphics[width=0.66\linewidth]{cap12/f_robustez}
\caption{Error de modelo: desde $t=0.5$\,s la masa es $0.12$\,kg. La ley exacta sola queda peor que la jacobiana; con integrador, el error se anula.}
\label{fig:fl-robustez}
\end{figure}

\codigorepo{cap12_linealizacion/C12_Ej9_LevitadorExacto.m}{robustez}{Exacta sola y exacta con integrador}

\begin{trampa}[label=tr:fl-geometria]{``Las propiedades geométricas son robustas''}
\plausible{El grado relativo y la forma normal son estructurales; un error pequeño de modelo no los cambia, así que el diseño basado en ellos es robusto.}
El grado relativo puede cambiar con perturbaciones pequeñas del modelo (un término despreciado que lleva la entrada a una derivada anterior de la salida), y la cancelación usa los valores numéricos de $a(x)$ y $b(x)$, no solo su estructura. Kokotović y Arcak lo resumen así: las propiedades geométricas pueden no ser robustas, y hay que combinarlas con herramientas analíticas~\citep{KoAr01}.
\correcto{Un diseño por linealización exacta se acompaña de un análisis de robustez (funciones de Lyapunov, ganancias, márgenes) y, para errores constantes, de acción integral.}
\end{trampa}

\begin{modelo}[LevitadorFL.slx]
\cmd{crearModelo\_LevitadorFL.m} arma el lazo con un bloque \texttt{MATLAB Function} para la ley exacta (llama a \cmd{leyExactaLevitador.m}), otro para la planta y un filtro $p^4/(s+p)^4$ que genera la referencia y sus tres derivadas a partir de un escalón de $0.05$ a $0.12$\,m. Simula, compara con \texttt{ode45} y deja \texttt{mPlanta} como parámetro para ensayar el error de modelo.
\end{modelo}

%% ---------------------------------------------------------------------
\section{Dinámica cero}
\label{sec:fl-cero}
%% ---------------------------------------------------------------------

\subsection{Un ejemplo que se ve perfecto}

\begin{ejemplo}[label=ej:fl-nominima]{La salida converge; el estado escapa}
Para
\begin{equation}
  \dot x_1=x_1+x_2,\qquad \dot x_2=3x_1^2x_2+x_1+u,\qquad y=-x_1^3+x_2,
  \label{eq:fl-nominima}
\end{equation}
diseñe una ley que lleve $y$ a cero exponencialmente y estudie qué pasa con $x$.
\solucion $L_gh=\jac hx\,(0,1)\T=1$: el grado relativo es $r=1$ en todo $\R^2$. Tomamos $\xi=h(x)=-x_1^3+x_2$ y $\eta=x_1$, que cumple $L_g\eta=0$; el jacobiano $\begin{bmatrix}-3x_1^2&1\\1&0\end{bmatrix}$ tiene determinante $-1$, y $x_2=\xi+\eta^3$. Entonces
\[
  \dot\xi=L_fh+u=\eta-3\eta^3+u,\qquad
  \dot\eta=x_1+x_2=\xi+\eta+\eta^3 .
\]
(La versión anterior de estas notas tenía $\dot\eta=\xi+3\eta^2+\eta^3$; lo correcto es lo de arriba.) Con $u=-\eta+3\eta^3+v$ y $v=-2\xi$, $\dot\xi=-2\xi$: la salida va a cero como $e^{-2t}$, como se diseñó. Pero $\eta$ obedece $\dot\eta=\xi+\eta+\eta^3$, y cuando $\xi\approx0$ queda $\dot\eta\approx\eta+\eta^3$, que escapa en tiempo finito. Desde $x(0)=(0.05,\ 0.6)$, la trayectoria llega a la curva $\{x_2=x_1^3\}$ (donde $y=0$), se desliza por ella y $\abs\eta$ alcanza $2.2$ en $t=1.394$, cuando $\abs y$ ya vale $0.037$ (\cref{fig:fl-nominima}); el control explota con ella.
\end{ejemplo}

\begin{figure}[ht]
\centering
\includegraphics[width=\linewidth]{cap12/f_fase_no_minima}
\caption{El sistema \eqref{eq:fl-nominima} con la ley linealizante. Izquierda: la trayectoria cae a $Z^*=\{x_2=x_1^3\}$ y se va por ella. Derecha: $y$ converge como se diseñó y $\eta=x_1$ escapa.}
\label{fig:fl-nominima}
\end{figure}

\video{V12_2_fase_no_minima}{La salida converge, el estado escapa}

\codigorepo{cap12_linealizacion/C12_Ej10_FaseNoMinima.m}{escape}{Linealización perfecta, estado que escapa}

La realimentación hizo el conjunto $\{x:h(x)=L_fh(x)=\dots=L_f^{r-1}h(x)=0\}$ \emph{invariante y atractivo}. Lo que pasa dentro de él decide la estabilidad.

\subsection{Anular la salida}

Sea $x^o$ un equilibrio de $\dot x=f(x)$ con $h(x^o)=0$ y grado relativo $r<n$; en forma normal, $z^o=(0,\eta^o)$. La pregunta es: ¿qué pares (condición inicial, entrada) producen $y(t)\equiv0$?

Si $y\equiv0$, también $\dot y\equiv\dots\equiv y^{(r-1)}\equiv0$: necesariamente $\xi(t)\equiv0$. Además, de \eqref{eq:fl-formanormal},
\[
  0=y^{(r)}=a(0,\eta)+b(0,\eta)\,u
  \quad\Longrightarrow\quad
  u^*(t)=-\frac{a\bigl(0,\eta(t)\bigr)}{b\bigl(0,\eta(t)\bigr)} ,
\]
que es la ley \eqref{eq:fl-ley} con $v=0$. Y lo que queda libre es $\eta$, que obedece
\begin{equation}
  \dot\eta=q(0,\eta),
  \label{eq:fl-cero}
\end{equation}
desde un $\eta(0)$ \emph{cualquiera}.

\begin{definicion}[label=def:fl-cero]{Dinámica cero y fase mínima \citep{ByIs88}}
La ecuación \eqref{eq:fl-cero} es la \clave{dinámica cero} del sistema \eqref{eq:fl-afin} con respecto a la salida $y=h(x)$: la dinámica residual cuando la salida se \emph{fuerza} a ser idénticamente cero. Si $\eta^o$ es un equilibrio asintóticamente estable de \eqref{eq:fl-cero}, el sistema es de \clave{fase mínima}; si es inestable, de \clave{fase no mínima}; si es estable pero no asintóticamente estable, de \clave{fase mínima débil}.
\end{definicion}

En el \cref{ej:fl-nominima}, la dinámica cero es $\dot\eta=\eta+\eta^3$: fase no mínima.

\subsection{En las coordenadas originales}

La condición $y^{(i-1)}(t)=0$, $i=1,\dots,r$, es $L_f^{i-1}h\bigl(x(t)\bigr)=0$. El estado tiene que quedarse en
\begin{equation}
  Z^*=\bigl\{x\in\R^n:\ h(x)=L_fh(x)=\dots=L_f^{r-1}h(x)=0\bigr\},
  \label{eq:fl-Zstar}
\end{equation}
que, por el \cref{lem:fl-independencia}, es una superficie (variedad) de dimensión $n-r$ cerca de $x^o$. La ley
\[
  u^*(x)=-\frac{L_f^rh(x)}{L_gL_f^{r-1}h(x)}
\]
hace $Z^*$ invariante, porque $\frac{d}{dt}L_f^{i-1}h=L_f^ih$ para $i<r$ y $\frac{d}{dt}L_f^{r-1}h=L_f^rh+L_gL_f^{r-1}h\,u^*=0$: en $Z^*$, todas las derivadas de las funciones que lo definen son cero.

\begin{proposicion}[label=prop:fl-ceroorig]{La dinámica cero sin coordenadas}
La dinámica cero es la restricción del campo $f^*(x)=f(x)+g(x)\,u^*(x)$ a la variedad $Z^*$. En particular, no depende de cuáles funciones complementarias $\phi_{r+1},\dots,\phi_n$ se escojan.
\end{proposicion}

\begin{trampa}[label=tr:fl-intrinseca]{``La dinámica cero no es una propiedad del sistema''}
\plausible{La dinámica cero aparece porque escogimos una entrada particular, $u^*$; es una consecuencia del controlador, no del sistema.}
Así lo decía la versión anterior de estas notas, y es engañoso. Para una terna $(f,g,h)$ dada, la dinámica cero \emph{no} depende de las coordenadas (\cref{prop:fl-ceroorig}) ni de la realimentación: si se aplica $u=\alpha(x)+\beta(x)v$ con $\beta\neq0$ y después se anula la salida, se llega a la misma variedad y a la misma dinámica. No se puede cambiar con control. Lo que sí la cambia es \emph{la salida}: van der Pol con $y=x_1$ no tiene dinámica cero ($r=n$); con $y=x_1^2+x_2$ ($r=1$, $\xi=x_1^2+x_2$, $\eta=x_1$) se tiene $\dot\eta=x_2=\xi-\eta^2$, y la dinámica cero es $\dot\eta=-\eta^2$, cuyo origen es estable solo desde un lado: no es de fase mínima.
\correcto{La dinámica cero es una propiedad del par planta--salida, como los ceros de una función de transferencia. Kokotović y Arcak lo subrayan: ni el grado relativo ni la dinámica cero se pueden modificar con realimentación~\citep{KoAr01}.}
\end{trampa}

\subsection{El caso lineal: la dinámica cero son los ceros}

\begin{ejemplo}[label=ej:fl-cerolineal]{Ceros de la función de transferencia}
Para $\dot x_1=x_2$, $\dot x_2=\alpha x_3+u$, $\dot x_3=\beta x_3-u$, $y=x_1$, calcule la dinámica cero y compárela con la función de transferencia.
\solucion $CB=0$ y $CAB=1$: $r=2$. Con $\eta=x_2+x_3$ ($L_g\eta=1-1=0$) y $x_3=\eta-\xi_2$,
\[
  \dot\eta=\alpha x_3+\beta x_3=(\alpha+\beta)(\eta-\xi_2),
  \qquad\text{dinámica cero: }\ \dot\eta=(\alpha+\beta)\,\eta .
\]
Y $G(s)=\dfrac{s-(\alpha+\beta)}{s^2(s-\beta)}$: el \emph{cero} está en $s=\alpha+\beta$. El sistema es de fase mínima si y solo si $\alpha+\beta<0$, es decir, si el cero está en el semiplano izquierdo. Con la ley $u=-\alpha x_3+v$, $v=-4y-4\dot y$ y $x(0)=(1,0,0)$, la salida es \emph{idéntica} con $(\alpha,\beta)=(1,-2)$ y $(2,-1)$; con el cero en $s=+1$, el estado interno y el control crecen como $e^{t}$: $x_3(6)=179.3$ contra $-0.0096$ (\cref{fig:fl-ceros}).
\end{ejemplo}

\begin{figure}[ht]
\centering
\includegraphics[width=\linewidth]{cap12/f_ceros}
\caption{La misma $y(t)$ con el cero en $s=-1$ (azul) y en $s=+1$ (rojo, discontinua). La diferencia está en lo que no se ve: $x_3$ y $u$ (escala logarítmica simétrica).}
\label{fig:fl-ceros}
\end{figure}

\codigorepo{cap12_linealizacion/C12_Ej11_DinamicaCero.m}{lineal}{La dinámica cero son los ceros}

\begin{ideaclave}
En un sistema lineal, la dinámica cero tiene como valores propios los ceros de $G(s)$. ``Fase mínima'' viene de ahí: con los ceros en el semiplano izquierdo, la fase de $G(\jw)$ es la mínima posible entre las funciones con la misma magnitud. La definición no lineal conserva el nombre y la consecuencia: lo que la salida no ve tiene que ser estable por sí mismo.
\end{ideaclave}

\begin{ejemplo}[label=ej:fl-cero3d]{Una superficie de dimensión dos}
Con $n=3$ y $r=1$, $Z^*$ es una superficie. Construyamos el sistema en forma normal, ya con la ley linealizante y $v=-3\xi$:
\[
  \dot\xi=-3\xi,\qquad \dot\eta=\begin{bmatrix}s&-2\\2&s\end{bmatrix}\eta+\begin{bmatrix}\xi\\0\end{bmatrix},
\]
y volvamos a coordenadas $x=\bigl(\eta_1,\ \eta_2,\ \xi+0.3(\eta_1^2-\eta_2^2)\bigr)$, de modo que $Z^*=\{x_3=0.3(x_1^2-x_2^2)\}$ es una silla de montar. Toda trayectoria cae a $Z^*$ como $e^{-3t}$, y después \emph{solo} la dinámica cero decide: con $s=-0.25$ el estado gira hacia el origen ($\norm{\eta(14)}=0.039$); con $s=+0.12$, se aleja girando ($\norm{\eta(14)}=6.86$), aunque $\xi$ sea la misma en los dos casos (\cref{fig:fl-cero3d}).
\end{ejemplo}

\begin{figure}[ht]
\centering
\includegraphics[width=\linewidth]{cap12/f_dinamica_cero3d}
\caption{La variedad $Z^*$ y una trayectoria que cae a ella. Con la dinámica cero estable (izquierda) llega al origen; con la inestable (derecha) se aleja por la superficie.}
\label{fig:fl-cero3d}
\end{figure}

\video{V12_3_dinamica_cero}{Después de llegar a $Z^*$, solo manda la dinámica cero}

\begin{trampa}[label=tr:fl-salida0]{``Si $y\to0$, el sistema está estabilizado''}
\plausible{La ley lleva la salida a cero exponencialmente con los polos escogidos; el problema de control está resuelto.}
La ley hace $Z^*$ atractivo, no el equilibrio. En un sistema de fase no mínima, el estado llega a $Z^*$ y se va por ella (\cref{ej:fl-nominima}): $y\to0$ mientras $x$ y $u$ explotan.
\correcto{Antes de usar la ley \eqref{eq:fl-ley} hay que calcular la dinámica cero. Si es inestable, la linealización entrada--salida no sirve para estabilizar con esa salida; hay que escoger otra salida o usar otra herramienta.}
\end{trampa}

%% ---------------------------------------------------------------------
\section{Estabilidad: de local a global}
\label{sec:fl-estabilidad}
%% ---------------------------------------------------------------------

\subsection{Fase mínima basta, localmente}

\begin{teorema}[label=thm:fl-local]{Estabilización local \citep{ByIs88}}
Suponga que el sistema \eqref{eq:fl-afin} tiene grado relativo $r$ en el equilibrio $x^o$, con $h(x^o)=0$, y que $\eta^o$ es un equilibrio \emph{localmente exponencialmente estable} de la dinámica cero \eqref{eq:fl-cero}. Entonces la ley
\[
  u=\frac{1}{b(z)}\bigl(-a(z)+v\bigr),\qquad v=-k_1\xi_1-k_2\xi_2-\dots-k_r\xi_r,
\]
con $s^r+k_rs^{r-1}+\dots+k_1$ Hurwitz, hace $x^o$ localmente exponencialmente estable.
\end{teorema}

\begin{proof}
En coordenadas $(\xi,\eta)$, el lazo cerrado es $\dot\xi=(A_c-B_cK)\xi$, $\dot\eta=q(\xi,\eta)$. Su jacobiano en $(0,\eta^o)$ es triangular por bloques,
\[
  J=\begin{bmatrix}A_c-B_cK&0\\[2pt] \dfrac{\partial q}{\partial\xi}(0,\eta^o)&\dfrac{\partial q}{\partial\eta}(0,\eta^o)\end{bmatrix},
\]
y sus valores propios son los de $A_c-B_cK$ (Hurwitz por diseño) y los de $\frac{\partial q}{\partial\eta}(0,\eta^o)$ (Hurwitz, porque la estabilidad exponencial de la dinámica cero equivale a que su linealización sea Hurwitz). Por el método indirecto de Lyapunov, $(0,\eta^o)$ es localmente exponencialmente estable, y como $\Phi$ es un difeomorfismo local, también lo es $x^o$.
\end{proof}

\begin{cuidado}[Asintótica no es exponencial]
Si la dinámica cero es asintóticamente estable pero no exponencialmente (por ejemplo $\dot\eta=-\eta^3$), $\frac{\partial q}{\partial\eta}(0,\eta^o)$ tiene un valor propio en el eje imaginario, y el método indirecto no concluye nada. La estabilidad local se puede probar todavía, pero con una función de Lyapunov (o con el teorema de la variedad central); no con la jacobiana.
\end{cuidado}

\begin{ejemplo}[label=ej:fl-faseminima]{Un sistema de fase mínima}
Estabilice el origen del sistema \eqref{eq:fl-ejkhalil}.
\solucion Con la forma normal del \cref{ej:fl-formakhalil}, anular la salida es $z_1=z_2=0$, es decir
\[
  Z^*=\{x\in\R^3:\ x_1=0,\ x_3=x_2^3\},
\]
y la dinámica cero es $\dot z_3=-z_3$: globalmente exponencialmente estable. El sistema es de fase mínima. Como $b=1+3x_2^2$ nunca se anula, la ley
\[
  u=\frac{v-(3x_2^3+x_1^2-x_3)}{1+3x_2^2},\qquad v=-2x_1-3(x_3-x_2^3),
\]
(polos en $-1$ y $-2$) está definida en todo $\R^3$, y por el \cref{thm:fl-local} el origen es localmente exponencialmente estable. En $Z^*$, $z_3=x_2+x_3\to0$ y $x_3=x_2^3$, así que $x_2,x_3\to0$. Desde $x(0)=(1,\ 0.5,\ -0.5)$, $\norm{x(10)}<10^{-3}$.
\end{ejemplo}

\codigorepo{cap12_linealizacion/C12_Ej11_DinamicaCero.m}{faseminima}{Un sistema de fase mínima con la ley linealizante}

\subsection{Global pide ISS}

El \cref{thm:fl-local} es local, aunque la dinámica cero sea \emph{globalmente} exponencialmente estable. Para una conclusión global hay que controlar cómo $\xi$ empuja a $\eta$ mientras converge: el lazo cerrado es una cascada, y la herramienta es la estabilidad entrada--estado.

\begin{proposicion}[label=prop:fl-global]{Estabilización global \citep[Lema~4.7]{Kha02}}
Si en coordenadas globales el subsistema $\dot\eta=q(\xi,\eta)$ es ISS con $\xi$ como entrada, la ley del \cref{thm:fl-local} hace el origen globalmente asintóticamente estable.
\end{proposicion}

La prueba es el resultado de cascadas de la unidad de ISS: $\xi(t)\to0$ exponencialmente, y un sistema ISS con una entrada que tiende a cero tiene su estado acotado y convergente a cero. Sin ISS, la conclusión falla aun en sistemas muy sencillos.

\begin{trampa}[label=tr:fl-global]{``Dinámica cero globalmente estable, estabilización global''}
\plausible{En $\dot\xi=-k\xi$, $\dot\eta=-\eta+\eta^2\xi$ la dinámica cero es $\dot\eta=-\eta$ (globalmente exponencialmente estable) y $\xi\to0$ exponencialmente; luego todo converge.}
El origen es localmente exponencialmente estable (\cref{thm:fl-local}), pero no global. Con $w=1/\eta$, $\dot w=w-\xi_0e^{-kt}$, y
\[
  w(t)=e^{t}\Bigl(w_0-\frac{\xi_0}{k+1}\Bigr)+\frac{\xi_0}{k+1}\,e^{-kt} .
\]
Si $\eta_0\xi_0>k+1$, $w$ se anula en tiempo finito: $\eta$ \emph{escapa a infinito}, en $t_e=\frac{1}{k+1}\ln\dfrac{\xi_0/(k+1)}{\xi_0/(k+1)-1/\eta_0}$. Con $k=1$ y $\xi_0=\eta_0=2$, $t_e=\tfrac12\ln2=0.3466$. Subir $k$ agranda la región que converge, pero nunca la vuelve global (\cref{fig:fl-escape}). Es la misma cascada de la \cref{tr:pa-cascada}: el subsistema $\eta$ no es ISS respecto a $\xi$, porque el término $\eta^2\xi$ crece más rápido de lo que $-\eta$ puede contener.
\correcto{La estabilidad global de la dinámica cero da estabilización local. Para la global hace falta ISS de $\dot\eta=q(\xi,\eta)$ respecto a $\xi$ (\cref{prop:fl-global}), o una ley que tenga en cuenta el acoplamiento (backstepping, pasivación: la \cref{thm:pa-cascada}).}
\end{trampa}

\begin{figure}[ht]
\centering
\includegraphics[width=\linewidth]{cap12/f_escape}
\caption{Trayectorias de $\dot\xi=-k\xi$, $\dot\eta=-\eta+\eta^2\xi$ en el plano $(\eta,\xi)$. Las que empiezan en la región sombreada ($\eta_0\xi_0>k+1$) escapan en tiempo finito (rojas). Con $k=5$ la región es más pequeña, pero no desaparece.}
\label{fig:fl-escape}
\end{figure}

\codigorepo{cap12_linealizacion/C12_Ej12_Global.m}{iss}{Escape en tiempo finito con dinámica cero estable}

\subsection{No cancelar lo que ayuda}

\begin{trampa}[label=tr:fl-cancelar]{``Linealizar por retroalimentación es cancelar todas las no linealidades''}
\plausible{Para $\dot x=x-x^3+u$, la ley natural es cancelar: $u=x^3-2x$ da $\dot x=-x$, un sistema lineal y estable.}
El término $-x^3$ \emph{ya} estabiliza cuando $\abs x$ es grande. Cancelarlo es pagar por quitar algo que ayudaba. Compare con $u=-2x$, que da $\dot x=-x-x^3$ (también GAS, con $V=\tfrac12x^2$). Desde $x(0)=3$: la cancelación pide $u(0)=21$ (contra $-6$); gasta $\int_0^\infty u^2\,\dif t=117/2=58.5$ (contra $2\ln10=4.6$, unas $12.7$ veces más), y converge más lento: $x(1)=3/e=1.104$ contra $0.372$ (\cref{fig:fl-nocancelar}).
\correcto{La linealización por retroalimentación es una herramienta de análisis y de diseño, no una obligación de cancelar todo. Los diseños ``constructivos'' (funciones de Lyapunov de control, backstepping, pasivación) conservan las no linealidades útiles~\citep{KoAr01}.}
\end{trampa}

\begin{figure}[ht]
\centering
\includegraphics[width=\linewidth]{cap12/f_no_cancelar}
\caption{$\dot x=x-x^3+u$ desde $x(0)=3$: cancelar (rojo) contra conservar el término estabilizante (azul). Estado, control y energía de control acumulada.}
\label{fig:fl-nocancelar}
\end{figure}

\codigorepo{cap12_linealizacion/C12_Ej12_Global.m}{nocancelar}{Cancelar contra conservar}

\subsection{El puente con la pasividad}

El \cref{cap:pasividad} terminó con una pregunta: ¿cuándo se puede volver pasivo un sistema con una realimentación preliminar? La respuesta está escrita en el lenguaje de este capítulo.

\begin{teorema}[label=thm:fl-biw]{Pasividad y fase mínima \citep{BIW91}}
El sistema \eqref{eq:fl-afin} es equivalente por retroalimentación (con $u=\alpha(x)+\beta(x)v$, $\beta\neq0$) a un sistema pasivo con almacenamiento definido positivo si y solo si tiene grado relativo uno en $x=0$ y es de fase mínima débil.
\end{teorema}

Las dos condiciones tienen sentido a la luz del \cref{cap:pasividad}. Si $\Vdot\le u\,y$ con $y=h=L_gV$, la entrada aparece en $\dot y$ a través de $L_gh=L_g(L_gV)$, que en el origen vale $g\T\,\nabla^2V\,g$: positiva si el hessiano de $V$ es definido positivo y $g(0)\neq0$, de modo que $r=1$. Y en $Z^*=\{y=0\}$, $\Vdot\le0$: la dinámica cero tiene a $V$ como función de Lyapunov, y es al menos débilmente estable. Con la equivalencia, la ley $v=-\phi(y)$ estabiliza por inyección de amortiguamiento (\cref{thm:pa-amortiguamiento}). Las dos últimas unidades del curso son dos caras de la misma estructura: la \emph{energía} (pasividad) y la \emph{geometría} (grado relativo y dinámica cero).

\begin{profundiza}[¿Cuándo existe una salida con $r=n$?]
Si existe una salida $h$ con grado relativo $n$, el sistema es linealizable entrada--estado y no hay dinámica cero. Jakubczyk y Respondek~\citep{JaRe80} y Hunt, Su y Meyer~\citep{HSM83} probaron que esa $h$ existe cerca de $x^o$ si y solo si
\begin{enumerate}
  \item la matriz $\bigl[\,g\ \ \mathrm{ad}_fg\ \ \cdots\ \ \mathrm{ad}_f^{n-1}g\,\bigr](x^o)$ tiene rango $n$ (el análogo de la controlabilidad), y
  \item la distribución $\operatorname{span}\{g,\mathrm{ad}_fg,\dots,\mathrm{ad}_f^{n-2}g\}$ es \emph{involutiva} cerca de $x^o$ (cerrada bajo el corchete de Lie),
\end{enumerate}
con $\mathrm{ad}_fg=[f,g]=\jac gxf-\jac fxg$ \citep[\S13.2]{Kha02}, \citep[\S4.2]{Isi95}. Aun si existe, $h$ puede no tener sentido físico: la salida que linealiza no tiene por qué ser la que se quiere regular. (Fuera del alcance de la evaluación del curso.)
\end{profundiza}

%% ---------------------------------------------------------------------
\section{Para llevar}
\label{sec:fl-cierre}
%% ---------------------------------------------------------------------

\begin{parallevar}
\begin{enumerate}
  \item \textbf{Jacobiana:} $u=u^*-K(x-\eq)$ da estabilidad \emph{local}; el modelo lineal cambia con el punto de operación. La acción integral corrige errores constantes si \eqref{eq:fl-rango} se cumple (no hay cero en $s=0$).
  \item \textbf{Derivada de Lie:} $L_f\lambda=\jac\lambda xf$, la derivada a lo largo de las soluciones ($\Vdot=L_fV$). \textbf{Difeomorfismo:} cambio de coordenadas no lineal; jacobiano no singular da uno \emph{local}, no uno global.
  \item \textbf{Grado relativo} $r$: cuántas veces hay que derivar $y$ para ver $u$. Depende de la salida y puede no existir en un punto ($b(x^o)=0$). En el caso lineal, es el exceso de polos sobre ceros.
  \item \textbf{Forma normal:} $\xi=(h,\dots,L_f^{r-1}h)$ es una cadena de $r$ integradores con $\dot\xi_r=a+bu$; $\dot\eta=q(\xi,\eta)$ no ve la entrada.
  \item \textbf{Linealización entrada--salida:} $u=(v-L_f^rh)/L_gL_f^{r-1}h$ deja $y^{(r)}=v$; un lazo externo lineal estabiliza o sigue referencias con $r$ derivadas. Si $r=n$, linealización entrada--estado.
  \item \textbf{Dinámica cero:} $\dot\eta=q(0,\eta)$, lo que queda cuando la salida se fuerza a cero; la restricción de $f+gu^*$ a $Z^*$. Depende de la salida, no de las coordenadas ni de la realimentación. En el caso lineal, sus valores propios son los ceros.
  \item \textbf{Estabilidad:} fase mínima con dinámica cero exponencialmente estable da estabilización local; para global hace falta ISS de $\eta$ respecto a $\xi$. Grado relativo uno y fase mínima débil equivalen a poder volver pasivo el sistema.
\end{enumerate}
\end{parallevar}

\begin{trampa}[title={\faBug\hspace{0.45em}Las trampas del capítulo, en una caja}]{}
\begin{itemize}
  \item Diseñar los polos no fija la respuesta lejos del punto de operación (\cref{tr:fl-local}); con integrador no quiere decir robusto (\cref{tr:fl-integral}).
  \item Exacta no es perfecta: la cancelación usa el modelo (\cref{tr:fl-robusta}); las propiedades geométricas pueden no ser robustas (\cref{tr:fl-geometria}).
  \item Jacobiano no singular en todas partes no da un difeomorfismo global (\cref{tr:fl-difeo}).
  \item El grado relativo puede no existir en un punto (\cref{tr:fl-nodefinido}) y depende de la salida (\cref{tr:fl-salida}).
  \item Un escalón no es una referencia admisible (\cref{tr:fl-escalon}); la ley exacta tiene singularidades con sentido físico (\cref{tr:fl-singular}).
  \item $y\to0$ no implica $x\to0$ (\cref{tr:fl-salida0}); la dinámica cero no la cambia el control, la cambia la salida (\cref{tr:fl-intrinseca}).
  \item Dinámica cero globalmente estable no da estabilización global (\cref{tr:fl-global}); no hay que cancelar lo que ayuda (\cref{tr:fl-cancelar}).
\end{itemize}
\end{trampa}

\begin{profundiza}[Correcciones respecto a las versiones anteriores (Handout9\_2 y las diapositivas de 2025)]
(i) En el ejemplo de fase no mínima, $\dot\eta=\xi+\eta+\eta^3$; la versión anterior tenía $\xi+3\eta^2+\eta^3$. (ii) La transformación de la salida es $\bar h(z)=h(\Phi^{-1}(z))$, no $\bigl[\jac\Phi xh(x)\bigr]_{x=\Phi^{-1}(z)}$. (iii) ``Las \emph{zero dynamics} no son una propiedad intrínseca del sistema'' era engañoso: no dependen de las coordenadas ni de la realimentación; sí dependen de la salida (\cref{tr:fl-intrinseca}). (iv) En el ejemplo con $\phi_3=1+x_1-e^{x_2}$, la ecuación de $\eta$ lleva un signo menos: $\dot\eta=-(-1+\eta+e^{\xi_2})(1+\xi_2e^{\xi_2})$. (v) En el ejemplo de fase mínima, la entrada que anula la salida es $u=-a/b$, no $a/b$. (vi) La forma normal de van der Pol con $y=x_1^2+x_2$ tenía erratas ($2z_z^3$, la ecuación de $\dot x_2$ en vez de $\dot z_2$); la correcta está en el \cref{ej:fl-coordvdp}. (vii) Van der Pol se escribe como en el \cref{cap:orbitas}, $\mu(1-x_1^2)x_2$, y no $(1-\epsilon x_1^2)x_2$. (viii) Masa del levitador con una sola letra, $m$. (ix) Se agregaron la condición de rango del control integral, la prueba del teorema local, la necesidad de ISS para un resultado global, las trampas de diseño de Kokotović y Arcak, y el puente con la pasividad.
\end{profundiza}

%% ---------------------------------------------------------------------
\section{Problemas}
\label{sec:fl-problemas}
%% ---------------------------------------------------------------------

Los problemas marcados con \etiquetaM{} requieren \matlab/\simulink; los marcados con \etiquetaEstrella{} son de profundización. Cada uno tiene una trampa: búsquela antes de dar la respuesta.

\subsection*{Herramientas y grado relativo}

\begin{problema}[un grado relativo que se pierde]\label{prob:fl-grado}
Para $\dot x_1=-x_1+x_2^2$, $\dot x_2=x_3$, $\dot x_3=u$, $y=x_1$: (a) calcule $L_gh$, $L_gL_fh$ y $L_gL_f^2h$; (b) ¿cuál es el grado relativo en un punto con $x_2\neq0$? (c) ¿Y en el origen? (d) ¿Se puede aplicar la ley \eqref{eq:fl-ley} para estabilizar el origen?
\end{problema}

\begin{problema}[un cero en el semiplano derecho]\label{prob:fl-markov}
Para $G(s)=\dfrac{1-s}{(s+1)^3}$: (a) calcule el grado relativo con los parámetros de Markov de una realización; (b) ¿qué signo tiene la respuesta al escalón para $t$ pequeño? ¿Y su valor final? (c) \etiquetaM{} Grafíquela. ¿Qué tiene que ver esto con la dinámica cero?
\end{problema}

\begin{problema}[¿difeomorfismos globales?]\label{prob:fl-difeo}
Decida cuáles son difeomorfismos globales de $\R^2$: (a) $(x_1+x_2^3,\ x_2)$; (b) $(\sinh x_1,\ x_2+x_1)$; (c) $(x_1^3,\ x_2)$; (d) $(e^{x_1},\ x_2)$.
\end{problema}

\subsection*{Forma normal, dinámica cero y estabilidad}

\begin{problema}[dónde se pone el sensor]\label{prob:fl-sensor}
Para el doble integrador $\dot x_1=x_2$, $\dot x_2=u$ con la salida $y=x_1+c\,x_2$: (a) ¿cuál es el grado relativo si $c\neq0$? ¿Y si $c=0$? (b) Con $c\neq0$, calcule la dinámica cero. (c) ¿Para qué $c$ es de fase mínima? Compare con el cero de $G(s)$.
\end{problema}

\begin{problema}[asintótica pero no exponencial]\label{prob:fl-iss}
Para $\dot x_1=-x_1^3+x_2$, $\dot x_2=x_1^2+u$, $y=x_2$: (a) calcule el grado relativo y la forma normal; (b) calcule la dinámica cero; (c) ¿se puede aplicar el \cref{thm:fl-local}? (d) Con $u=-x_1^2-k\,x_2$, $k>0$, ¿es el origen globalmente asintóticamente estable? (Pista: ISS.)
\end{problema}

\begin{problema}[van der Pol con otra salida]\label{prob:fl-vdpcero}
Para van der Pol con $y=x_1^2+x_2$ (\cref{ej:fl-coordvdp}), suponga que se aplica la entrada $u^*$ que mantiene $y\equiv0$ desde una condición inicial en $Z^*$. (a) Resuelva la dinámica cero $\dot\eta=-\eta^2$ desde $\eta_0<0$. (b) ¿Qué le pasa a $x$? (c) ¿Es de fase mínima débil?
\end{problema}

\begin{problema}[global, gracias a ISS]\label{prob:fl-expiss}
\etiquetaEstrella{} Para el sistema del \cref{ej:fl-formaexp}: (a) muestre que $1+\xi_2e^{\xi_2}\ge1-1/e$ para todo $\xi_2$; (b) concluya que $\dot\eta=-(-1+\eta+e^{\xi_2})(1+\xi_2e^{\xi_2})$ es ISS con entrada $\xi$; (c) diseñe una ley que estabilice globalmente el origen.
\end{problema}

\subsection*{Aplicaciones}

\begin{problema}[cuán rápido se puede bajar la esfera]\label{prob:fl-filtro}
\etiquetaM{} En el levitador, la referencia es un escalón de amplitud $\Delta$ pasado por $p^4/(s+p)^4$. (a) Muestre que la aceleración de la esfera cumple $\ddot y\le g-\frac km\dot y$: no puede caer más rápido que en caída libre. (b) Muestre que el máximo de $\ddot y_d$ es $c_4\,\Delta\,p^2$, con $c_4=\max_{\tau\ge0}\tfrac16(3\tau^2-\tau^3)e^{-\tau}$, y calcule $c_4$. (c) Estime el mayor $p$ para $\Delta=0.15$\,m y compruébelo con \cmd{C12\_Ej9\_LevitadorExacto.m}.
\end{problema}

\begin{problema}[seguimiento del péndulo]\label{prob:fl-pendulo}
\etiquetaM{} Para el péndulo \eqref{eq:fl-pendulo} con $y=x_1$: (a) diseñe la ley de seguimiento de $y_d=\tfrac\pi2+\tfrac\pi4\sin t$ con polos en $-2$, $-2$; (b) simúlela con $\hat a=a$ y con $\hat a=8$; (c) agregue un integrador del error y repita. ¿Desaparece el error con $\hat a=8$? ¿Por qué no del todo?
\end{problema}

\begin{problema}[cancelar o no]\label{prob:fl-cancelar}
\etiquetaEstrella{} Para $\dot x_1=x_2$, $\dot x_2=-x_1^3-x_2+u$: (a) diseñe una ley por linealización exacta que lleve el origen a $\ddot x_1+2\dot x_1+x_1=0$; (b) muestre que $u=-x_1-x_2$, que conserva el resorte $-x_1^3$, estabiliza globalmente el origen (busque una función de energía); (c) compare las dos leyes desde $x(0)=(3,0)$: control máximo y convergencia. ¿Qué pierde cada una?
\end{problema}

%% ---------------------------------------------------------------------
\section{Soluciones seleccionadas}
%% ---------------------------------------------------------------------

\paragraph{Problema~\ref{prob:fl-grado}}
(a) $L_gh=0$; $L_fh=-x_1+x_2^2$, $L_gL_fh=0$; $L_f^2h=x_1-x_2^2+2x_2x_3$ y $L_gL_f^2h=2x_2$. (b) $r=3$ donde $x_2\neq0$. (c) En el origen, $L_gL_f^2h(0)=0$ pero $L_gL_f^2h$ no es idénticamente nula cerca de $0$: no hay grado relativo (\cref{tr:fl-nodefinido}). (d) No: la ley divide por $2x_2$, que se anula justo en el equilibrio que se quiere estabilizar.

\paragraph{Problema~\ref{prob:fl-markov}}
(a) $CB=0$, $CAB=-1$: $r=2$. (b) Cerca de $t=0$, $y\approx CAB\,t^2/2=-t^2/2<0$, pero $y(\infty)=G(0)=1>0$: la respuesta empieza en la dirección equivocada (respuesta inversa). (c) El cero en $s=+1$ es el valor propio de la dinámica cero: el sistema es de fase no mínima, y la respuesta inversa es su huella en el tiempo.

\paragraph{Problema~\ref{prob:fl-difeo}}
(a) Sí: inversa $(z_1-z_2^3,\ z_2)$, suave. (b) Sí: $\sinh$ es biyectiva con inversa suave, y $x_2=z_2-\operatorname{arsinh}z_1$. (c) No: es biyectiva, pero su inversa $z_1^{1/3}$ no es diferenciable en $0$, y $\det\jac\Phi x=3x_1^2$ se anula en $x_1=0$. (d) No: no es sobreyectiva ($\Phi(\R^2)$ es el semiplano $z_1>0$); es un difeomorfismo de $\R^2$ en su imagen, pero no de $\R^2$ en $\R^2$.

\paragraph{Problema~\ref{prob:fl-sensor}}
(a) $L_gh=c$: $r=1$ si $c\neq0$ y $r=2$ si $c=0$. (b) Con $\xi=x_1+cx_2$ y $\eta=x_1$ ($L_g\eta=0$), $\dot\eta=x_2=(\xi-\eta)/c$, y la dinámica cero es $\dot\eta=-\eta/c$. (c) Fase mínima si y solo si $c>0$. $G(s)=(1+cs)/s^2$ tiene su cero en $s=-1/c$: en el semiplano izquierdo si y solo si $c>0$. Medir la posición \emph{adelantada} ($c>0$) da un sistema de fase mínima; medirla \emph{atrasada}, uno de fase no mínima.

\paragraph{Problema~\ref{prob:fl-iss}}
(a) $L_gh=1$: $r=1$. Con $\xi=x_2$ y $\eta=x_1$: $\dot\xi=\eta^2+u$, $\dot\eta=-\eta^3+\xi$. (b) $\dot\eta=-\eta^3$: asintóticamente (no exponencialmente) estable. (c) No: el teorema pide estabilidad exponencial (\cref{thm:fl-local} y el cuidado que lo sigue). (d) Sí. El lazo es $\dot\xi=-k\xi$, $\dot\eta=-\eta^3+\xi$, y $\dot\eta=-\eta^3+\xi$ es ISS respecto a $\xi$: con $V=\tfrac12\eta^2$, $\Vdot=-\eta^4+\eta\xi\le-\tfrac12\eta^4$ si $\abs\eta\ge(2\abs\xi)^{1/3}$. Por la \cref{prop:fl-global}, el origen es GAS. El teorema local no aplicaba; el resultado global sí.

\paragraph{Problema~\ref{prob:fl-vdpcero}}
(a) $\eta(t)=\eta_0/(1+\eta_0t)$, que escapa a $-\infty$ en $t=-1/\eta_0$. (b) $x_1=\eta\to-\infty$ y $x_2=-x_1^2\to-\infty$ en tiempo finito, con $y\equiv0$. (c) No: el origen de $\dot\eta=-\eta^2$ no es estable (desde $\eta_0<0$, por pequeño que sea, la solución se aleja).

\paragraph{Problema~\ref{prob:fl-filtro}}
(a) $\ddot y=\xi_3=g-\frac km\dot y-\frac{L_0a\,i^2}{2m(a+y)^2}$, y el último término es no negativo. (b) La respuesta al escalón de $1/(s+1)^4$ tiene derivada $\tau^3e^{-\tau}/6$ y segunda derivada $(3\tau^2-\tau^3)e^{-\tau}/6$, máxima en $\tau=3-\sqrt3$: $c_4=0.1306$. Con $\tau=pt$, $\ddot y_d=\Delta p^2\cdot(\dots)$. (c) $0.1306\cdot0.15\cdot p^2<g$ da $p<22.4$. En la simulación, con $p=22$ la corriente baja hasta $1.2$\,A y se recupera; con $p=23$ se anula.

\paragraph{Problema~\ref{prob:fl-cancelar}}
(a) $u=x_1^3+x_2+v$ con $v=-x_1-2x_2$, es decir $u=x_1^3-x_1-x_2$. (b) $u=-x_1-x_2$ da $\dot x_2=-x_1^3-x_1-2x_2$; con $V=\tfrac14x_1^4+\tfrac12x_1^2+\tfrac12x_2^2$ (agregue $\tfrac12x_1^2$ por el resorte lineal que introdujo la ley), $\Vdot=-2x_2^2\le0$, y LaSalle da GAS. (c) La cancelación pide $u(0)=24$ para \emph{quitar} un resorte que empujaba hacia el origen; la otra nunca pide más de $3$ en valor absoluto y llega antes cerca del origen ($\norm{x(2)}=0.74$ contra $1.46$), aunque al principio el resorte la lanza con más velocidad. Lo que pierde la segunda es la respuesta lineal exacta, la que se diseñó.

\respuestaspausas

